%        1         2         3         4         5         6         7         8         9         0         1         2
%23456789012345678901234567890123456789012345678901234567890123456789012345678901234567890123456789012345678901234567890
\documentclass[12pt]{MIA-USA}

\usepackage{layout}
\usepackage{subcaption}
\usepackage{paralist}
\usepackage{breakcites}
\usepackage{natbib}
% Comillas tipogr\'aficas con \enquote{}: el car\'acter " es un atajo activo de babel en espa\~nol.
\usepackage[autostyle]{csquotes}

\usepackage[pagebackref=true,breaklinks=true]{hyperref}
\hypersetup{
	pdftitle={Evaluaci\'on por componentes de un pipeline de Retrieval-Augmented Generation para la b\'usqueda de empleo remoto: el caso JobAI},
	pdfauthor={Luis Esteban Pe\~na Rodr\'iguez}, 
	pdfsubject={}, 
	pdfcreator={DarMart with Latex}, 
	colorlinks=true,
	linkcolor=GrisUno,
	citecolor=AzulInstitucional,
	filecolor=AzulClaro,
	urlcolor=AzulInstitucional,
	linktoc=all
}

\graphicspath{{Images/}}

% babel (espa\~nol) vuelve activos < y >; esta librer\'ia evita el conflicto con las flechas de TikZ.
\usetikzlibrary{babel}
% Numeraci\'on progresiva cap\'itulo.secci\'on (NTC 1486), para que las referencias
% cruzadas entre cap\'itulos no sean ambiguas.
\renewcommand{\thesection}{\thechapter.\arabic{section}}
% NTC 1486 usa "Tabla"; babel en espa\~nol usa "Cuadro" por defecto.
\addto\captionsspanish{\renewcommand{\tablename}{Tabla}}

% Marca visible para datos que se completar\'an con los resultados de la evaluaci\'on.
\newcommand{\pendiente}[1]{\textcolor{red}{[Pendiente: #1]}}


\title{Evaluaci\'on por componentes de un pipeline de \textit{Retrieval-Augmented Generation} para la b\'usqueda de empleo remoto: el caso JobAI}
\author{Luis Esteban Pe\~na Rodr\'iguez}
\documenttype{Trabajo de grado}

\advisor{Joaqu\'in Fernando S\'anchez Cifuentes}

% Si hay un coasesor 
% \coadvisor{Nombres y apellidos Co-Asesor}


% Aqui comienza el documento
% - -- --- ----- ------- ------- ----- --- -- -
\begin{document}
    % \layout{}
	
	\maketitle
	% P\'aginas opcionales de plantilla: se completan en la versi\'on final.
	% % Página de aceptación. Contiene las firmas del presidente o director y de los jurados que participan en la revisión, sustentación y aprobación del trabajo. Adicionalmente, incluye la ciudad y la fecha de entrega (día, mes, año), conservando los márgenes establecidos

{\hfill
    \begin{minipage}{0.5\textwidth}
        \centering
        
        Nota de aceptaci\'on:\\
        \vspace*{0.8cm}
        
        \hrulefill\\[0.2cm]
        \hrulefill\\[0.2cm]
        \hrulefill\\[0.2cm]
        \hrulefill\\[0.2cm]
        \hrulefill\\[0.2cm]
        \hrulefill\\
        \vspace*{4cm}
        
        \hrulefill\\
        Firma del presidente del jurado\\
        \vspace*{2.5cm}
        
        \hrulefill\\
        Firma del jurado\\
        \vspace*{2.5cm}
        
        \hrulefill\\
        Firma del jurado\\
        \vspace*{2.5cm}
    
    \end{minipage}
}
\newpage
	% %  Nota mediante la cual el autor ofrece su trabajo, en forma especial, a personas o entidades. Su presentación es opcional y debe conservar los márgenes
\subsection*{Dedicatoria}
\vspace*{1.5cm}
{\hfill
    \begin{minipage}{0.5\textwidth}
        \noindent
         Nota mediante la cual el autor ofrece su trabajo, en forma especial, a personas o entidades~\cite{NTC14862008}. Su presentación es opcional y, aunque cada autor podr\'{a} determinar la distribuci\'{o}n del texto en la p\'{a}gina, se sugiere esta presentaci\'{o}n. En ella el autor dedica su trabajo en forma especial a personas y/o entidades.\\
         \vspace*{0.5cm}
         
         Por ejemplo:\\
         \\
         A mis padres\\
         \vspace*{1.0cm}
         
         o\\
         \vspace*{1.0cm}
         
        A computer would deserve to be called intelligent if it could deceive a human into believing that it was human.\\\\
        Alan Turing\\
    
    \end{minipage}
}
\newpage
	% % Página de agradecimientos. En ella el (los) autor (es) expresa (n) el reconocimiento hacia las personas y entidades que asesoraron técnicamente, suministraron datos, financiaron total o parcialmente la investigación o contribuyeron significativamente al desarrollo del tema. Es opcional y debe contener, además de la nota correspondiente, los nombres de las personas con sus respectivos cargos y los nombres completos de las instituciones y su aporte al trabajo.

\section*{Agradecimientos}

Página de agradecimientos. En ella el (los) autor (es) expresa (n) el reconocimiento hacia las personas y entidades que asesoraron técnicamente, suministraron datos, financiaron total o parcialmente la investigación o contribuyeron significativamente al desarrollo del tema. Es opcional y debe contener, además de la nota correspondiente, los nombres de las personas con sus respectivos cargos y los nombres completos de las instituciones y su aporte al trabajo~\cite{NTC14862008}.


%  Se debe dejar la nueva página para guardar el orden de las seciones
\newpage
	% Resumen
 	% Presentación abreviada y precisa del contenido de un documento, sin agregar interpretación o crítica; no debe exceder una página.
\section*{Resumen}

\pendiente{Cifras provisionales: juez de relevancia local, conjunto de consultas en borrador y validación humana del juez pendiente. Se actualizan con la corrida final.}

La búsqueda de empleo remoto se reparte entre varias plataformas cuyos motores comparan las palabras de la consulta con las de la oferta, y fallan cuando ambas describen el mismo cargo con términos distintos. Este trabajo desarrolla y evalúa JobAI, una plataforma de búsqueda de empleo remoto basada en un pipeline de generación aumentada por recuperación (RAG), con el objetivo de superar las limitaciones de precisión de la búsqueda por palabras clave. Siguiendo el paradigma de Design Science Research, el trabajo constituye una nueva iteración del ciclo de diseño-evaluación del artefacto. En la fase de diseño se corrigió la consulta vectorial, que no usaba el índice, y se migró el almacenamiento vectorial de pgvector a Qdrant detrás de una interfaz común. La evaluación se realizó sobre un corpus congelado de 1.068 ofertas reales (4.013 fragmentos) y 100 consultas: 75 en inglés, en cinco categorías, y 25 en español. La recuperación densa, servida por pgvector y por Qdrant, se comparó con la búsqueda por subcadena que ofrece la plataforma (ILIKE) y con BM25. Los juicios de relevancia los asignó un modelo de lenguaje sobre un agrupamiento de profundidad 10. La fidelidad de las respuestas se midió con RAGAS, la latencia por etapa del pipeline y el compromiso entre \textit{recall} y latencia de ambos motores sobre el corpus real y conjuntos sintéticos de 100.000 y 500.000 vectores, con y sin filtros.

La recuperación densa obtuvo un nDCG@10 de 0,73, frente a 0,56 de BM25 y 0,06 de ILIKE. Las diferencias fueron significativas en las cinco métricas según la prueba de aleatorización pareada con corrección de Holm. La mayor ventaja sobre BM25 se dio en las consultas con desajuste de vocabulario (+0,21). En español, el nDCG@10 de la recuperación densa estuvo entre 0,64 y 0,66, y el de BM25 fue de 0,34. La búsqueda aproximada con HNSW no perdió calidad frente a la búsqueda exacta. En las consultas de agregación, el clasificador eligió la ruta y la plantilla SQL esperadas en 12 de 15 casos, y la respuesta coincidió con la consulta de referencia en 9. Con el juez local, la fidelidad media fue de 0,663 y la relevancia de la respuesta de 0,599. Sin embargo, la revisión de los casos con puntaje nulo mostró errores del juez, por lo que estas cifras no se consideran una medida fiable. Servido desde Qdrant, el contexto se construyó en 59~ms en la mediana, frente a 80,5~ms con pgvector, cuyo planificador recorre la tabla completa a esta escala. La generación dominó la latencia: el primer \textit{token} llegó en unos 0,53~s y la respuesta completa en unos 1,21~s. En el benchmark sobre el corpus real y sin filtro, pgvector igualó o superó el \textit{recall} de Qdrant con menor latencia hasta $ef = 128$. Con un filtro del 50\,\% y sin barrido iterativo, pgvector devolvió los 10 resultados solo en el 25,7\,\% de las consultas, mientras que Qdrant mantuvo un \textit{recall} de 1,000 con todos los filtros. \pendiente{Resultados con 100.000 y 500.000 vectores.}

\textbf{Palabras clave}: generación aumentada por recuperación, búsqueda semántica, búsqueda de empleo, bases de datos vectoriales, pgvector, Qdrant, evaluación de la recuperación de información, Design Science Research.

\newpage
\section*{Abstract}

\pendiente{English version of the abstract.}

\textbf{Keywords}: \pendiente{keywords}.

	
    
    % Crea la tabla de contenidos a partir de la estructura del documento
    \tableofcontents
    % Crea la lista de figuras (Opcional)
    \listoffigures
    % Crea la lista de tablas (Opcional)
    \listoftables
    % Lista de s\'imbolos y glosario: se completan en la versi\'on final.
    % \chapter*{Lista de s\'{\i}mbolos}
% \addcontentsline{toc}{chapter}{\numberline{}Lista de s\'{\i}mbolos}
\addcontentsline{toc}{chapter}{\hspace{1em} Lista de s\'{\i}mbolos}
Esta secci\'{o}n es opcional, dado que existen disciplinas que no manejan s\'{\i}mbolos y/o abreviaturas. A continuaci\'on se presentan algunos ejemplos tomados de la plantilla de Tesis de Maestr\'ia y doctorado de la Universidad Nacional.\\

Se incluyen s\'{\i}mbolos generales (con letras latinas y griegas), sub\'{\i}ndices, super\'{\i}ndices y abreviaturas (incluir s\'{o}lo las clases de s\'{\i}mbolos que se utilicen). Cada una de estas listas debe estar ubicada en orden alfab\'{e}tico de acuerdo con la primera letra del s\'{\i}mbolo.
\section*{S\'{\i}mbolos con letras latinas}
 \label{simbolos}
 \renewcommand{\arraystretch}{1.3}
%\begin{longtable}[l]{*{4}{>{$}l<{$}}p{9cm}}
\begin{longtable}[l]{>{$}l<{$}l>{$}l<{$}>{$}l<{$}}
%\begin{tabular}
\textbf{S\'{\i}mbolo}&\textbf{T\'{e}rmino}&\textbf{Unidad SI}&\textbf{Definici\'{o}n}\\[0.5ex]\hline
\endfirsthead%
\textbf{S\'{\i}mbolo}&\textbf{T\'{e}rmino}&\textbf{Unidad SI}&\textbf{Definici\'{o}n}\\[0.5ex]\hline
\endhead%
      A              &\'{A}rea                                   &\text{m}^{2}                         &\int\int dxdy\\%
      A_{\text{BET}} &\'{A}rea interna del s\'{o}lido                &\frac{\text{m}^{2}}{\text{g}}        &\text{ver DIN ISO 9277}\\%
      A_{\text{g}}   &\'{A}rea transversal de la fase gaseosa    &\text{m}^{2}                         &\text{Ec...}\\%
      A_{\text{s}}   &\'{A}rea transversal de la carga a granel  &\text{m}^{2}                         &\text{Ec...}\\%
      a              &Coeficiente                            &1                                    &\text{Ec...}\\%
      a              &Contenido de ceniza                    &1                                    &\frac{m_{\text{ceniza}}}{m_{\text{bm,0}}}\\%
      c              &Contenido de carbono                   &1                                    &\frac{m_{\text{C}}}{m}\\%
      c              &Longitud de la cuerda                  &\text{m}                             &\text{Figura...}\\
      c              &Concentraci\'{o}n de la cantidad de materia&\frac{\text{mol}}{\text{m}^{3}}      &\frac{n}{V}\\%
      D              &Di\'{a}metro                               &\text{m}                             &\\%
      E_{\text{A}}   &Energ\'{\i}a de activaci\'{o}n                  &\frac{\text{kJ}}{\text{mol}}         &\text{Ec....}\\%
      F              &Fracci\'{o}n de materia vol\'{a}til            &1                                    &\text{ver DIN 51720}\\%
      Fr             &N\'{u}mero de Froude                       &1                                    &\frac{\omega^{2}R}{g_{\text{0}}}\\%
      \overrightarrow{g}&Aceleraci\'{o}n de la gravedad          &\frac{\text{m}}{\text{s}^{2}}        &\frac{d^{2}\overrightarrow{r}}{dt^{2}}\\%
      H              &Entalp\'{\i}a                               &\text{J}                             &U+PV\\%
      H_{\text{o}}   &Poder calor\'{\i}fico superior              &\frac{\text{MJ}}{\text{kg}}          &\text{ver DIN 51857}\\%
      h              &Contenido de hidr\'{o}geno                 &1                                    &\frac{m_{\text{H}}}{m}\\%
      K              &Coeficiente de equilibrio              &1                                    &\text{Ec...}\\%
      L              &Longitud                               &\text{m}                             &DF\\%
      L              &Longitud del reactor                   &\text{m}                             &\text{Figura...}\\%
      m              &Masa                                   &\text{kg}                            &DF\\%
      \dot{m}        &Flujo de masa                          &\frac{\text{kg}}{\text{s}}           &\frac{m}{t}\\%
      n              &Velocidad de rotaci\'{o}n                  &\frac{\text{1}}{\text{s}}            &\frac{\omega}{2\pi}\\%
      n              &Cantidad de materia                    &\text{mol}                           &DF\\%
      P              &Presi\'{o}n                                &\text{Pa}                            &\frac{\vec{F}\cdot\vec{n}}{A}\\%
      Q              &Calor                                  &\text{kJ}                            &\text{1. $LT$}\\%
      T              &Temperatura                            &\text{K}                             &DF\\%
      t              &Tiempo                                 &\text{s}                             &DF\\%
      x_{\text{i}}   &Fracci\'{o}n de la cantidad de materia     &1                                    &\frac{n_{\text{i}}}{n}\\%
      V              &Volumen                                &\text{m}^{3}                         &\int{dr^{3}}\\%
      \vec{u}        &Velocidad                              &\frac{\text{m}}{\text{s}}            &(\frac{dr}{dt},r\frac{d\upsilon}{dt},\frac{dz}{dt})\\%
      w_{\text{i}}   &Fracci\'{o}n en masa del componente i      &1                                    &\frac{m_{\text{i}}}{m_{\text{0}}}\\%
      w_{\text{w,i}} &Contenido de humedad de la sustancia i &1                                    &\frac{m_{\text{\wasser}}}{m_{\text{i,0}}}\\%
      Z              &Factor de gases reales                 &1                                    &\frac{pv}{RT}\\%
\end{longtable}
\vspace{5ex}
\section*{S\'{\i}mbolos con letras griegas}

\begin{longtable}[l]{>{$}l<{$}l>{$}l<{$}>{$}l<{$}}
\textbf{S\'{\i}mbolo}&\textbf{T\'{e}rmino}&\textbf{Unidad SI}&\textbf{Definici\'{o}n}\\[0.5ex] \hline%
\endfirsthead%
\textbf{S\'{\i}mbolo}&\textbf{T\'{e}rmino}&\textbf{Unidad SI}&\textbf{Definici\'{o}n}\\[0.5ex] \hline%
\endhead%
\renewcommand{\arraystretch}{1.3}
 \label{simbolosg}
     \alpha_{\text{BET}}  &Factor de superficie                  &\frac{\text{m}^{2}}{\text{g}}   &(w_{\text{F,waf}})(A_{\text{BET}})\\%
     \beta_{\text{i}}     &Grado de formaci\'{o}n del componente i   &1                               &\frac{m_{\text{i}}}{m_{\text{bm,0}}}\\%
     \gamma               &Wandhaftreibwinkel (Stahlblech)       &1                               &\text{Secci\'{o}n...}\\
     \epsilon             &Porosidad de la part\'{\i}cula             &1                               &1-\frac{\rho_{\text{s}}}{\rho_{\text{w}}}\\%
     \eta                 &mittlere Bettneigungswinkel (St\"{u}rzen) &1                               &\text{Figura...}\\%
     \theta               &\'{A}ngulo de inclinaci\'{o}n de la cama      &1                               &\text{Figura...}\\
     \theta_{\text{O}}    &\'{A}ngulo superior de avalancha          &1                               &\text{Figura...}\\
     \theta_{\text{U}}    &\'{A}ngulo inferior de avalancha          &1                               &\text{Figura...}\\
     \kappa               &Velocidad de calentamientoe           &\frac{\text{K}}{\text{s}}       &\frac{dT}{dt}\\%
     \nu                  &Coeficiente estequiom\'{e}trico           &1                               &\text{ver DIN 13345}\\%
     \rho_{\text{b}}      &Densidad a granel                     &\frac{\text{kg}}{\text{m}^{3}}  &\frac{m_{\text{S}}}{V_{\text{S}}}\;(\text{Secci\'{o}n...})\\
     \rho_{\text{s}}      &Densidad aparente                     &\frac{\text{kg}}{\text{m}^{3}}  &\frac{m_{\text{F}}}{V_{\text{P}}}\;(\text{Secci\'{o}n...})\\
     \rho_{\text{w}}      &Densidad verdadera                    &\frac{\text{kg}}{\text{m}^{3}}  &\frac{m_{\text{F}}}{V_{\text{F}}}\;(\text{Secci\'{o}n...})\\
     \tau                 &Tiempo adimensional                   &1                               &\text{Ec....}\\%
     \Phi_{\text{V}}      &Flujo volum\'{e}trico                     &\frac{\text{m}^{3}}{\text{s}}   &\frac{\Delta V}{\Delta t}\\
     \omega               &Velocidad angular                     &\frac{1}{\text{s}}              &\frac{d\upsilon}{dt}\\

\end{longtable}


\section*{Sub\'{\i}ndices}
\begin{longtable}[l]{>{}l<{}l}
  \textbf{Sub\'{\i}ndice} & \textbf{T\'{e}rmino} \\[0.5ex] \hline%
  \endfirsthead%
  \textbf{Sub\'{\i}ndice} & \textbf{T\'{e}rmino} \\[0.5ex] \hline%
  \endhead%
\renewcommand{\arraystretch}{1.4}\label{simbolosg}

 bm&materia org\'{a}nica\\%
 DR&Dubinin-Radushkevich\\%
 E&Experimental\\%
 g&Fase gaseosa\\%
 k&Condensado\\%
 Ma&Macroporos\\%
 P&Part\'{\i}cula\\%
 p&Poro\\%
 p&Pirolizado\\%
 R&Reacci\'{o}n\\%
 t&Total\\%
 wf&Libre de agua\\%
 waf&Libre de agua y de ceniza\\%
 0&Estado de referencia\\%

\end{longtable}


\setlength{\extrarowheight}{0pt}


\section*{Super\'{\i}ndices}
\begin{longtable}[l]{>{}l<{}l}
  \textbf{Super\'{\i}ndice} & \textbf{T\'{e}rmino} \\[0.5ex] \hline%
  \endfirsthead%
  \textbf{Super\'{\i}ndice} & \textbf{T\'{e}rmino} \\[0.5ex] \hline%
  \endhead%
\renewcommand{\arraystretch}{1.4}\label{simbolosg}

 n &Coeficiente x\\%



\end{longtable}


\setlength{\extrarowheight}{0pt}


\section*{Abreviaturas}
\begin{longtable}[l]{>{}l<{}l}
  \textbf{Abreviatura} & \textbf{T\'{e}rmino} \\[0.5ex] \hline%
  \endfirsthead%
  \textbf{Abreviatura} & \textbf{T\'{e}rmino} \\[0.5ex] \hline%
  \endhead%
\renewcommand{\arraystretch}{1.4}\label{simbolosg}
 1.$LT$&Primera ley de la termodin\'{a}mica\\%
 $DF$    &Dimensi\'{o}n fundamental\\%
 $RFF$   &Racimos de fruta fresca\\%

\end{longtable}


\setlength{\extrarowheight}{0pt}
 	% % Lista alfabética de términos y sus definiciones o explicaciones necesarios para la comprensión del documento. La existencia de un glosario no justifica la omisión de una explicación en el texto la primera vez que aparece un término. El título glosario se escribe en mayúscula sostenida, centrado, a 3 cm del borde superior de la hoja.  El primer término aparece a dos interlíneas del título glosario, contra el margen izquierdo. Los términos se escriben con mayúscula sostenida seguidos de dos puntos y en orden alfabético. La definición correspondiente se coloca después de los dos puntos, se deja un espacio y se inicia con minúscula. Si ocupa más de un renglón, el segundo y los subsiguientes comienzan contra el margen izquierdo. Entre término y término se deja una interlínea. Su uso es opcional

\chapter*{Glosario}

Lista alfabética de términos y sus definiciones o explicaciones necesarios para la comprensión del documento. La existencia de un glosario no justifica la omisión de una explicación en el texto la primera vez que aparece un término. El título glosario se escribe en mayúscula sostenida, centrado, a 3 cm del borde superior de la hoja.  El primer término aparece a dos interlíneas del título glosario, contra el margen izquierdo. Los términos se escriben con mayúscula sostenida seguidos de dos puntos y en orden alfabético. La definición correspondiente se coloca después de los dos puntos, se deja un espacio y se inicia con minúscula. Si ocupa más de un renglón, el segundo y los subsiguientes comienzan contra el margen izquierdo. Entre término y término se deja una interlínea. Su uso es opcional~\cite{NTC14862008}
    
    
    % A partir de este punto se crea la estructura del documento con un archivo por capitulo conservando la siguiente jerarquía
    % chapter (Capitulo)
    %   - section (sección)
    %       - subsection (subsección)
    %           - subsubsection (subsubsección)
    
    % Introduction
    \chapter{Introducci\'on}
    % De acuerdo con la NTC 1486, la introducción presenta la importancia, el origen, los objetivos, el alcance, las limitaciones y la metodología del trabajo, sin anticipar resultados ni conclusiones (2 a 4 páginas).

Buscar empleo remoto exige hoy recorrer varias plataformas, cada una con su propio catálogo, sus propios filtros y un mecanismo de búsqueda que, en lo esencial, compara las palabras de la consulta con las palabras de la oferta. Ese mecanismo funciona cuando quien busca y quien publica describen el cargo con los mismos términos, y falla en cuanto no lo hacen. "Backend Developer" y "Software Engineer", "JS" y "JavaScript", "trabajo con datos" y "Analytics Engineer" pueden referirse a lo mismo sin compartir una sola palabra. La literatura de recuperación de información conoce este problema como \emph{desajuste de vocabulario}. La industria del reclutamiento a gran escala lo ha abordado desde hace casi una década con representaciones semánticas densas \citep{HaThuc2017,Shen2024}, pero las herramientas a disposición de quien busca empleo apenas lo reflejan.

La generación aumentada por recuperación (\textit{Retrieval-Augmented Generation}, RAG) \citep{Lewis2020} ofrece una alternativa. En lugar de exigir coincidencias exactas, recupera los fragmentos de información semánticamente más cercanos a la consulta y delega en un modelo de lenguaje la síntesis de una respuesta fundamentada en ellos. Aplicado a la búsqueda de empleo, este enfoque permite formular preguntas que un panel de filtros no puede expresar: qué ofertas de \textit{backend} en el propio \textit{stack} superan cierto salario, o cuáles de las ofertas guardadas se ajustan a la experiencia descrita en un currículo.

Este trabajo parte de JobAI, una plataforma de búsqueda de empleo remoto construida previamente por el autor, que agrega ofertas de tres fuentes públicas y expone un copiloto conversacional basado en un pipeline RAG. JobAI existe y funciona, pero nunca se había evaluado de forma sistemática. No se sabía, con evidencia cuantitativa, en qué medida su recuperación semántica supera a la búsqueda por palabra clave que pretende reemplazar. Tampoco se sabía qué tan fieles son sus respuestas a los datos recuperados, ni si la decisión inicial de almacenar los vectores en PostgreSQL mediante pgvector era adecuada frente a una base de datos vectorial dedicada. Esos tres vacíos dan origen a la investigación.

\section*{Objetivos}

El objetivo general es desarrollar y evaluar una plataforma de búsqueda de empleo remoto que utiliza un pipeline RAG para permitir la búsqueda semántica y conversacional sobre ofertas laborales, superando las limitaciones de precisión de la búsqueda por palabras clave. Para alcanzarlo se plantean cuatro objetivos específicos:
\begin{enumerate}
    \item Caracterizar, a partir de la literatura, las limitaciones de los sistemas actuales de búsqueda de empleo y las oportunidades que ofrece RAG.
    \item Diseñar y documentar la arquitectura del pipeline, incluidos la recuperación con filtros, la búsqueda híbrida, el reordenamiento y la migración del almacenamiento vectorial a una base de datos dedicada.
    \item Definir una metodología de evaluación con métricas cuantitativas.
    \item Evaluar el artefacto frente a un enfoque de búsqueda tradicional por palabras clave.
\end{enumerate}
El Capítulo~\ref{cap:problema} desarrolla el planteamiento completo.

\section*{Metodología}

El trabajo se enmarca en el paradigma de Design Science Research \citep{Hevner2004,Peffers2007}, que concibe la investigación como el diseño y la evaluación rigurosa de un artefacto que resuelve un problema relevante. Como el artefacto ya existía, se adopta el punto de entrada del proceso centrado en el diseño y desarrollo, y el trabajo se plantea como una nueva iteración del ciclo de diseño-evaluación. Esa iteración tiene dos partes. La de diseño consiste en corregir y migrar el almacenamiento vectorial y en revisar el artefacto frente a su documentación. La de evaluación incluye tres mediciones:
\begin{itemize}
    \item La calidad de la recuperación, frente a dos líneas base léxicas, sobre un corpus real congelado y un conjunto de consultas en cinco categorías.
    \item La fidelidad de las respuestas generadas, con las métricas de RAGAS \citep{Es2024}.
    \item La latencia del pipeline y un benchmark de pgvector frente a Qdrant, con datos reales y sintéticos, con y sin filtros estructurados.
\end{itemize}
Los juicios de relevancia los asigna un modelo de lenguaje, validado contra una muestra de juicios humanos.

\section*{Alcance y aporte}

El estudio se limita a JobAI y a su dominio: ofertas de empleo remoto en inglés, mayoritariamente del sector tecnológico, obtenidas de Remotive, Jobicy y Findwork. Evalúa la calidad técnica del pipeline y no la experiencia de usuarios reales, y compara dos representantes de los enfoques de almacenamiento vectorial, una extensión sobre un motor relacional y un motor dedicado, sin pretender una revisión exhaustiva del mercado.

El aporte es doble. En lo aplicado, ofrece una evaluación cuantitativa e integrada de un sistema RAG completo en producción en un dominio donde la literatura de RAG y la de emparejamiento de empleo rara vez se han cruzado \citep{RojasGaleano2022}. En lo metodológico, documenta cómo aplicar Design Science Research a un artefacto preexistente, incluido el hallazgo de que la propia preparación de la evaluación reveló defectos que la habrían invalidado.

\section*{Estructura del documento}

Además de esta introducción, el documento tiene seis capítulos:
\begin{itemize}
    \item El Capítulo~\ref{cap:problema} plantea el problema, los objetivos y el marco conceptual.
    \item El Capítulo~\ref{cap:estado-arte} revisa el estado del arte en seis líneas: RAG, emparejamiento de empleo, límites de la búsqueda léxica, Design Science Research, bases de datos vectoriales y clasificación híbrida de consultas.
    \item El Capítulo~\ref{cap:metodologia} detalla la metodología de evaluación.
    \item El Capítulo~\ref{cap:artefacto} documenta la arquitectura de JobAI y la iteración de diseño realizada.
    \item El Capítulo~\ref{cap:resultados} presenta y discute los resultados.
    \item El último capítulo responde a cada objetivo específico y propone trabajo futuro.
\end{itemize}

    
    % Capitulo 1
    \chapter{Planteamiento del problema y marco te\'orico}
    \label{cap:problema}
    \section{Planteamiento del problema}
\label{sec:problema}

La búsqueda de empleo remoto es hoy una experiencia fragmentada. Las ofertas se dispersan entre decenas de plataformas —agregadores generalistas, portales especializados y bolsas de empleo corporativas—, cada una con su propio esquema de filtros y su propio motor de búsqueda. La herramienta subyacente que casi todas comparten, sin embargo, ha cambiado poco en las últimas dos décadas: el emparejamiento por palabra clave (\textit{keyword-matching}) entre los términos de la consulta y los campos estructurados o el texto libre de la oferta.

Esta dependencia del emparejamiento léxico exacto es también la base de la mayoría de los sistemas de seguimiento de candidatos (\textit{Applicant Tracking Systems}, ATS) usados del lado del reclutador. La propia motivación de trabajos como el de \citet{HaThuc2017} —que documenta la migración de LinkedIn desde una búsqueda de talento basada en palabras clave (\textit{Query-By-Keyword}) hacia una basada en ejemplos (\textit{Query-By-Example})— y el de \citet{Bocharova2023} —que atribuye explícitamente a la variación léxica entre títulos de vacantes equivalentes ("Software Engineer" frente a "Backend Developer", por ejemplo) una pérdida sistemática de coincidencias relevantes— evidencia que el problema no es anecdótico sino estructural: el emparejamiento exacto falla precisamente en los casos en que existe sinonimia entre habilidades, variación en los títulos de cargo o ambigüedad en la formulación de la consulta. La revisión bibliométrica de \citet{RojasGaleano2022} sobre investigación en inteligencia artificial aplicada al emparejamiento empleo-currículo confirma que esta limitación ha sido, durante más de una década, uno de los motivadores centrales de la investigación en la intersección de la recuperación de información y el reclutamiento asistido por cómputo.

Frente a este panorama, la generación aumentada por recuperación (\textit{Retrieval-Augmented Generation}, RAG), propuesta originalmente por \citet{Lewis2020}, ofrece una alternativa arquitectónica distinta: en lugar de exigir una coincidencia léxica exacta, el sistema recupera los fragmentos de información semánticamente más cercanos a la consulta —representados como vectores en un espacio de \textit{embeddings}— y delega en un modelo de lenguaje la síntesis de una respuesta fundamentada en dichos fragmentos. Como documenta la revisión de \citet{Gao2023}, este paradigma ha probado ser particularmente efectivo para reducir la alucinación de los modelos de lenguaje y para responder preguntas que dependen de una base de conocimiento externa y actualizable, dos propiedades directamente aplicables al dominio de la búsqueda de empleo: el conocimiento (las ofertas vigentes) cambia constantemente, y una respuesta no fundamentada en datos reales (por ejemplo, un salario o un requisito inventado) es inaceptable para quien la usa para tomar una decisión de carrera.

La mayoría de las plataformas de empleo que hoy incorporan "inteligencia artificial" lo hacen de forma superficial: un asistente conversacional genérico superpuesto a los mismos filtros de siempre, o un puntaje de similitud aislado que no reemplaza al motor de búsqueda. JobAI —el artefacto sobre el que se desarrolla este trabajo (véase el Capítulo 4)— se construyó bajo una premisa distinta: que el copiloto de IA \emph{sea} el motor de búsqueda, mediante un pipeline RAG completo que recupera, filtra y sintetiza información sobre ofertas laborales agregadas de múltiples fuentes, en lugar de limitarse a conversar sobre datos que el usuario ya podría ver en un panel de filtros. Sin embargo, aunque JobAI ya existe como artefacto funcional, dos vacíos persisten y motivan el presente trabajo. Primero, un vacío de evaluación: no existe todavía una medición cuantitativa, sistemática, de qué tan bien cumple el pipeline RAG su promesa de precisión y fidelidad frente al enfoque de búsqueda por palabra clave que busca superar. Segundo, un vacío de fundamentación arquitectónica: la decisión de almacenamiento vectorial —inicialmente resuelta con la extensión pgvector sobre PostgreSQL por su simplicidad operativa— no ha sido contrastada empíricamente frente a una base de datos vectorial dedicada, a pesar de que la literatura reciente en sistemas de gestión de bases de datos vectoriales (revisada en el Capítulo 2) documenta de manera consistente las limitaciones de escalabilidad de la primera aproximación.

\section{Formulación del problema}

De lo anterior se deriva la pregunta de investigación que orienta este trabajo:

\begin{quote}
\textit{¿En qué medida un pipeline de Retrieval-Augmented Generation, implementado sobre una plataforma real de búsqueda de empleo remoto y evaluado mediante métricas cuantitativas de precisión de recuperación, fidelidad de las respuestas y latencia, permite superar las limitaciones del emparejamiento por palabra clave, y cómo debe diseñarse dicha evaluación para que sus resultados sean válidos y reproducibles?}
\end{quote}

Esta pregunta tiene una segunda dimensión, subordinada pero necesaria para responderla con rigor metodológico: dado que JobAI ya es un artefacto construido, ¿cómo se enmarca metodológicamente la evaluación de un artefacto preexistente, y qué papel juega dentro de ese marco la decisión de migrar el almacenamiento vectorial de pgvector a una base de datos vectorial dedicada (Qdrant)? El Capítulo 3 retoma esta segunda dimensión en detalle, apoyándose en el marco de Design Science Research introducido conceptualmente en la Sección~\ref{sec:marco-teorico-dsr} de este capítulo y revisado críticamente en el Capítulo 2.

\section{Objetivos}

\subsection{Objetivo general}

Desarrollar y evaluar una plataforma de búsqueda de empleo remoto que utiliza un pipeline de Retrieval-Augmented Generation para permitir búsqueda semántica y conversacional sobre ofertas laborales, superando las limitaciones de precisión de la búsqueda por palabras clave.

\subsection{Objetivos específicos}

\begin{enumerate}
    \item Caracterizar, a partir de la literatura, las limitaciones de los sistemas actuales de búsqueda de empleo y las oportunidades que ofrece RAG para superarlas.
    \item Diseñar y documentar la arquitectura del pipeline RAG (ingestión, \textit{chunking}, \textit{embeddings}, clasificación de consultas, síntesis) aplicada al dominio de búsqueda de empleo, incluyendo la migración del almacenamiento vectorial de pgvector a una base de datos vectorial dedicada.
    \item Definir una metodología de evaluación con métricas cuantitativas (precisión de recuperación, fidelidad/\textit{groundedness} de las respuestas, tiempo de respuesta) para validar la plataforma.
    \item Evaluar los resultados obtenidos y contrastarlos frente a un enfoque de búsqueda tradicional por palabras clave.
\end{enumerate}

\section{Justificación}

La justificación de este trabajo se sostiene sobre tres ejes complementarios.

Desde el punto de vista \textbf{práctico}, la fragmentación y las limitaciones de precisión descritas en la Sección~\ref{sec:problema} tienen un costo directo para quien busca empleo: tiempo perdido evaluando ofertas irrelevantes y oportunidades genuinamente compatibles que nunca aparecen porque el sistema de búsqueda no reconoció la equivalencia semántica entre lo que la persona escribió y lo que la oferta describe. Una plataforma que resuelve esto con evidencia cuantitativa de su desempeño —y no solo con la promesa genérica de "búsqueda con IA"— tiene valor inmediato y verificable.

Desde el punto de vista \textbf{teórico}, la revisión bibliométrica de \citet{RojasGaleano2022} muestra que la investigación en inteligencia artificial aplicada al emparejamiento empleo-currículo es un campo activo pero todavía concentrado en tareas de \textit{ranking} y recomendación (como los sistemas de \citealp{HaThuc2017,Shen2024,Bocharova2023} revisados en el Capítulo 2); la aplicación específica de arquitecturas RAG completas —con recuperación, síntesis conversacional y clasificación híbrida de consultas— al dominio de búsqueda de empleo es todavía un espacio de trabajo emergente. Este trabajo aporta, en ese sentido, un caso de estudio evaluado empíricamente en un subdominio donde la literatura de RAG general \citep{Lewis2020,Gao2023} y la literatura específica de matching de empleo rara vez se han integrado explícitamente.

Desde el punto de vista \textbf{metodológico}, este trabajo constituye una aplicación del paradigma de Design Science Research \citep{Hevner2004,Peffers2007} a un artefacto de software ya construido y en producción, en lugar de a un artefacto diseñado desde cero para fines exclusivamente académicos. Esto exige adaptar el ciclo de diseño-evaluación de Peffers et al. a un escenario de iteración sobre un sistema existente —la migración del almacenamiento vectorial y el diseño de la evaluación cuantitativa son, en los términos de ese marco, una nueva iteración del ciclo—, lo cual se documenta en detalle en el Capítulo 3 y puede sensiblemente extenderse a un artefacto de similar naturaleza en trabajos de investigación aplicada.

\section{Alcance y limitaciones}

El alcance de este trabajo se circunscribe a JobAI como artefacto de estudio: una plataforma que agrega ofertas de empleo remoto de tres fuentes públicas (Remotive, Jobicy y Findwork), las indexa mediante \textit{embeddings} vectoriales y expone un copiloto conversacional basado en un pipeline RAG (véase el Capítulo 4 para el detalle arquitectónico). La evaluación cuantitativa (Capítulo 5) se limita a las métricas de precisión de recuperación, fidelidad de las respuestas y latencia definidas en el Capítulo 3, contrastadas contra un \textit{baseline} de búsqueda por palabras clave sobre el mismo conjunto de datos.

Quedan fuera del alcance: (i) una evaluación de experiencia de usuario o de adopción con usuarios reales, dado que el enfoque es la calidad técnica del pipeline y no su recepción; (ii) la generalización de los resultados a otros dominios de búsqueda distintos al de empleo remoto; (iii) la comparación exhaustiva contra la totalidad de bases de datos vectoriales disponibles en el mercado, restringiéndose la comparación empírica a pgvector y Qdrant como representativas de los dos enfoques arquitectónicos discutidos en el Capítulo 2 (extensión vectorial sobre un motor relacional frente a motor vectorial dedicado); y (iv) la evaluación de la calidad de los datos de origen (Remotive, Jobicy, Findwork), que se asumen como dados.

\section{Marco teórico conceptual}

Esta sección introduce, a nivel conceptual, el vocabulario técnico necesario para seguir el resto del documento. La revisión crítica de cada una de estas líneas —con sus fuentes primarias, sus hallazgos específicos y su relación explícita con las decisiones de diseño de JobAI— se desarrolla en el Capítulo 2.

\subsection{Recuperación de información y el problema del emparejamiento léxico}

La recuperación de información clásica enfrenta un problema bien documentado: el \textit{desajuste de vocabulario} (\textit{vocabulary mismatch}), en el que dos textos semánticamente equivalentes no comparten los términos exactos necesarios para que un sistema de coincidencia léxica los relacione. En el dominio de reclutamiento, este problema se manifiesta como sinonimia entre habilidades (p. ej., "JS" y "JavaScript"), variación entre títulos de cargo equivalentes, y ambigüedad en la forma en que una persona describe lo que busca frente a cómo una oferta describe lo que ofrece. Los sistemas de búsqueda por palabra clave y los ATS tradicionales, al depender de la coincidencia exacta (o de expansión limitada mediante sinónimos predefinidos), heredan esta limitación de forma directa.

\subsection{Modelos de lenguaje de gran escala}

Los modelos de lenguaje de gran escala (\textit{Large Language Models}, LLMs) son modelos generativos preentrenados sobre grandes volúmenes de texto, capaces de producir lenguaje natural fluido y de razonar sobre instrucciones formuladas en lenguaje natural. Su limitación central, señalada de forma recurrente en la literatura de RAG \citep{Gao2023}, es que su conocimiento está codificado paramétricamente en los pesos del modelo en el momento del entrenamiento: no tienen acceso nativo a información posterior a su corte de entrenamiento ni a datos privados o específicos de un dominio (como el catálogo de ofertas de una plataforma de empleo), y tienden a "alucinar" —generar afirmaciones plausibles pero no sustentadas— cuando se les pide responder sobre información que no poseen.

\subsection{Retrieval-Augmented Generation (RAG)}

RAG, propuesto por \citet{Lewis2020}, combina un componente de \textit{recuperación} (que busca los fragmentos de una base de conocimiento externa más relevantes para una consulta dada, típicamente mediante similitud vectorial) con un componente \textit{generativo} (un LLM que sintetiza una respuesta a partir de la consulta y los fragmentos recuperados). La ventaja central de este diseño es que ancla la respuesta del modelo en información verificable y actualizable, en lugar de depender exclusivamente de su conocimiento paramétrico. La encuesta de \citet{Gao2023} organiza la evolución de este paradigma en tres niveles de sofisticación creciente —RAG ingenuo (recuperación y generación en un solo paso), RAG avanzado (con reformulación de consultas, reordenamiento de resultados y filtrado de contexto) y RAG modular (con componentes intercambiables y flujos condicionales, como enrutamiento entre distintas estrategias de recuperación)—, una taxonomía que resulta directamente pertinente para ubicar el diseño del pipeline de JobAI (Capítulo 4). Evaluar un sistema RAG plantea, a su vez, un reto propio: al no existir siempre una respuesta de referencia contra la cual comparar, \citet{Es2024} proponen con RAGAS un conjunto de métricas de evaluación sin referencia humana —fidelidad (\textit{faithfulness}, si la respuesta se sostiene en los fragmentos recuperados), relevancia de la respuesta, y precisión/exhaustividad del contexto recuperado— que este trabajo retoma como base conceptual para el diseño de las métricas de evaluación del Capítulo 3.

\subsection{Bases de datos vectoriales y búsqueda por similitud}

Un \textit{embedding} es una representación vectorial de un fragmento de texto (o de cualquier otro dato) en un espacio de alta dimensión, construida de forma que la distancia (típicamente similitud coseno) entre dos vectores refleje la similitud semántica entre los textos que representan. Buscar los fragmentos más relevantes para una consulta se convierte, en este espacio, en un problema de búsqueda del vecino más cercano (\textit{nearest neighbor search}). Dado que la búsqueda exacta de vecinos más cercanos es computacionalmente costosa a gran escala, en la práctica se recurre a algoritmos de búsqueda \textit{aproximada} del vecino más cercano (ANN); el más influyente de ellos es HNSW \citep{Malkov2020}, una estructura de grafo jerárquico de navegación que permite búsquedas en tiempo logarítmico. Estos algoritmos pueden implementarse tanto como una extensión sobre un motor de base de datos relacional (como pgvector sobre PostgreSQL) como dentro de un motor construido específicamente para datos vectoriales (como Qdrant o Milvus). El Capítulo 2 revisa en detalle la evidencia empírica que sustenta la decisión de migrar de la primera aproximación a la segunda.

\subsection{Design Science Research como paradigma metodológico}
\label{sec:marco-teorico-dsr}

Design Science Research (DSR) es un paradigma de investigación en sistemas de información orientado a la creación y evaluación rigurosa de artefactos —métodos, modelos, constructos o instanciaciones de software— que resuelven problemas organizacionales o técnicos relevantes. \citet{Hevner2004} formalizan este paradigma mediante siete lineamientos (entre ellos, que el artefacto debe ser evaluado rigurosamente y que la investigación debe producir una contribución verificable), mientras que \citet{Peffers2007} lo operacionalizan en un proceso de seis actividades secuenciales —identificación del problema, definición de objetivos, diseño y desarrollo, demostración, evaluación y comunicación— que admite entradas en distintos puntos del ciclo, incluyendo la iteración sobre un artefacto ya existente. Este trabajo adopta DSR como marco metodológico general; su aplicación concreta a JobAI —incluyendo cómo se interpreta la migración a Qdrant como una nueva iteración del ciclo diseño-evaluación— se desarrolla en el Capítulo 3.

    
    % Capitulo 2
    \chapter{Estado del arte}
    \label{cap:estado-arte}
    Este capítulo revisa críticamente la literatura correspondiente a las seis líneas identificadas en la Fase 1 del proyecto de tesis: (1) el fundamento técnico de RAG, (2) los sistemas de recomendación y emparejamiento de empleo basados en NLP y LLMs, (3) los límites del emparejamiento por palabra clave y de los ATS en reclutamiento, (4) Design Science Research como paradigma metodológico, (5) las bases de datos vectoriales, y (6) la clasificación híbrida de consultas y el text-to-SQL. Cada sección cierra con un párrafo que vincula explícitamente los hallazgos revisados con una decisión de diseño concreta de JobAI.

\section{Fundamento técnico de RAG}

El paradigma de \textit{Retrieval-Augmented Generation} fue introducido por \citet{Lewis2020} como una arquitectura que combina un recuperador denso (\textit{Dense Passage Retriever}) con un modelo generativo tipo secuencia-a-secuencia (BART), entrenados de forma conjunta para producir respuestas fundamentadas en pasajes recuperados de una base de conocimiento no paramétrica —en su caso, Wikipedia. Los autores demuestran que este diseño supera a los modelos puramente paramétricos en tareas intensivas en conocimiento (\textit{knowledge-intensive NLP tasks}), y que la separación entre memoria paramétrica (el modelo) y memoria no paramétrica (el índice de recuperación) permite actualizar el conocimiento del sistema sin reentrenar el modelo generativo, simplemente actualizando el índice.

La encuesta de \citet{Gao2023} sistematiza la evolución posterior de este paradigma en tres categorías: RAG \textit{ingenuo}, que recupera y genera en un único paso sin refinamiento adicional; RAG \textit{avanzado}, que incorpora preprocesamiento de la consulta (reformulación, expansión), reordenamiento de los resultados recuperados y compresión o filtrado del contexto antes de pasarlo al modelo generativo; y RAG \textit{modular}, que descompone el pipeline en componentes intercambiables y permite flujos condicionales, como enrutar la consulta hacia distintas estrategias de recuperación según su naturaleza. Los autores identifican como retos abiertos del paradigma la calidad de la recuperación (fragmentos irrelevantes degradan la respuesta final incluso con un modelo generativo capaz), la gestión de la longitud de contexto, y la persistencia de alucinaciones cuando el contexto recuperado es insuficiente o contradictorio.

Evaluar un sistema RAG es, en sí mismo, un problema no trivial: a diferencia de tareas con una respuesta de referencia única, la calidad de una respuesta generada depende de si está \emph{fundamentada} en el contexto recuperado y de si el contexto recuperado era, a su vez, relevante para la consulta. \citet{Es2024} proponen RAGAS, un framework de evaluación sin referencia humana que opera a partir de tres dimensiones: fidelidad (\textit{faithfulness}, el grado en que cada afirmación de la respuesta puede inferirse del contexto recuperado), relevancia de la respuesta (\textit{answer relevancy}) y precisión/exhaustividad del contexto recuperado (\textit{context precision/recall}). Estas métricas son especialmente relevantes para este trabajo porque no requieren un conjunto de referencia anotado manualmente para cada consulta posible sobre un catálogo de ofertas laborales que cambia cada seis horas.

\textbf{Relación con JobAI.} El pipeline implementado en JobAI —descrito en detalle en el Capítulo 4— corresponde, en la taxonomía de \citet{Gao2023}, a un diseño de RAG modular: la consulta del usuario se aumenta con información del perfil (títulos y habilidades extraídos del currículo) antes de embeberse, el contexto recuperado se somete a una etapa de detección de campos que complementa fragmentos incompletos con datos estructurados (salario, beneficios, requisitos), y la consulta completa se enruta condicionalmente entre una ruta de recuperación semántica y una ruta de agregación SQL (véase la Sección~\ref{sec:estado-arte-hibrido}). El diseño de las métricas de evaluación del Capítulo 3 —precisión de recuperación y fidelidad de las respuestas— se apoya directamente en las dimensiones de fidelidad y precisión de contexto propuestas por RAGAS \citep{Es2024}, adaptadas al dominio de ofertas laborales.

\section{Sistemas de recomendación y emparejamiento de empleo con NLP y LLMs}

La industria del reclutamiento en línea ha sido, durante la última década, uno de los espacios donde la búsqueda semántica basada en representaciones densas se ha adoptado con mayor decisión, precisamente por la limitación estructural del emparejamiento léxico descrita en la Sección~\ref{sec:estado-arte-keyword}. \citet{HaThuc2017} documentan la transición de LinkedIn Talent Search desde un modelo de \textit{Query-By-Keyword} —donde el reclutador construye manualmente una consulta booleana— hacia uno de \textit{Query-By-Example}, donde el reclutador simplemente indica uno o varios candidatos ideales y el sistema aprende a generar la consulta y el modelo de \textit{ranking} a partir de esos ejemplos. Los autores destacan como reto principal la ausencia de datos de uso históricos para un paradigma completamente nuevo, resuelto mediante una estrategia de aprendizaje que reutiliza registros del sistema anterior.

En una línea similar mas centrada en la representación semántica de vacantes, \citet{Bocharova2023} proponen VacancySBERT, una red siamesa basada en \textit{Sentence-BERT} entrenada para representar títulos de vacantes y habilidades en un espacio vectorial donde la similitud coseno refleja equivalencia semántica, incluso cuando el título exacto difiere ("Software Engineer" y "Backend Developer", por ejemplo). Los autores reportan mejoras de hasta 21,5\,\% sobre \textit{baselines} como FastText, BERT y SentenceBERT genéricos al incorporar las habilidades mencionadas en la oferta como característica adicional durante la inferencia, evidencia directa de que una representación especializada al dominio de reclutamiento supera a una representación general de propósito amplio.

A una escala mayor, \citet{Shen2024} describen cómo LinkedIn incorporó técnicas de recuperación aprendida (\textit{learning to retrieve}) —incluyendo recuperación densa y recuperación generativa— a su sistema de búsqueda y recomendación de empleo, en sustitución parcial del paradigma clásico de dos etapas (indexación invertida por entidades extraídas más re-ranking), que los autores señalan como costoso de mantener manualmente y limitado frente a consultas que no calzan exactamente con los términos indexados.

Finalmente, la revisión bibliométrica de \citet{RojasGaleano2022} sitúa estos desarrollos dentro de una tendencia más amplia: analizando la producción científica sobre inteligencia artificial aplicada al emparejamiento empleo-currículo, los autores documentan un crecimiento sostenido del área en la última década, concentrado mayoritariamente en tareas de análisis de currículos, extracción de habilidades y \textit{ranking}/recomendación mediante representaciones semánticas, con una presencia todavía incipiente de arquitecturas conversacionales o generativas completas aplicadas específicamente a este dominio.

\textbf{Relación con JobAI.} La estrategia de aumentar la consulta del usuario con los títulos y habilidades extraídos de su currículo antes de embeberla —descrita en el Capítulo 4— sigue el mismo principio que motiva a \citet{HaThuc2017} y \citet{Bocharova2023}: representar semánticamente tanto la necesidad del usuario como la oferta, en lugar de exigir una coincidencia textual entre ambas. La brecha señalada por \citet{RojasGaleano2022} respecto a la escasez de arquitecturas conversacionales completas en este dominio es, precisamente, el espacio que este trabajo busca documentar y evaluar empíricamente.

\section{Límites del emparejamiento por palabra clave y de los ATS en reclutamiento}
\label{sec:estado-arte-keyword}

Más allá de que la literatura de matching semántico (Sección anterior) motive sus propuestas contrastándolas con el emparejamiento por palabra clave, vale la pena precisar en qué consiste exactamente esa limitación. El problema, conocido en recuperación de información como \textit{desajuste de vocabulario} (\textit{vocabulary mismatch}), ocurre porque un sistema de coincidencia léxica exacta —o de expansión mediante listas fijas de sinónimos— no puede anticipar todas las formas en que dos personas describirán el mismo concepto. \citet{HaThuc2017} lo ilustran desde la perspectiva del reclutador: traducir un conjunto complejo de criterios de contratación a una consulta booleana efectiva es una tarea que requiere pericia y que, aun así, omite candidatos cuyo perfil describe la misma experiencia con palabras distintas a las de la consulta. \citet{Bocharova2023} lo cuantifican desde la perspectiva del catálogo de vacantes: dos títulos de cargo pueden ser funcionalmente equivalentes y, sin embargo, no compartir ningún término, lo que un ATS basado en coincidencia exacta interpreta simplemente como "no encontrado".

Esta limitación se traduce en dos tipos de error complementarios: falsos negativos (una oferta relevante no aparece porque no comparte el vocabulario exacto de la consulta) y falsos positivos de baja calidad (una oferta aparece por compartir una palabra clave aislada, sin que el resto de su contenido sea relevante). Ninguno de los dos es evidente para quien realiza la búsqueda, que simplemente percibe un sistema con baja precisión, sin distinguir si el problema es de recuperación insuficiente o de ausencia de contexto para razonar sobre la relevancia.

\textbf{Relación con JobAI.} Esta limitación es, textualmente, el punto de partida del proyecto (véase la Sección~\ref{sec:problema}): la premisa de que un panel de filtros y una caja de búsqueda por palabra clave no pueden expresar consultas del tipo \textit{"qué roles de backend en mi stack pagan más de \$120k"} o \textit{"cuáles de mis ofertas guardadas calzan con mi experiencia en Python"}, que requieren razonar sobre significado y no solo sobre coincidencia textual.

\section{Bases de datos vectoriales: de una extensión relacional a un motor dedicado}
\label{sec:estado-arte-vectordb}

Un sistema RAG requiere, en su componente de recuperación, resolver eficientemente el problema de búsqueda del vecino más cercano sobre potencialmente millones de vectores de \textit{embedding}. \citet{Malkov2020} proponen HNSW (\textit{Hierarchical Navigable Small World graphs}), un algoritmo de búsqueda aproximada del vecino más cercano basado en un grafo de proximidad organizado en capas jerárquicas, que permite una complejidad de búsqueda logarítmica y que se ha convertido en el índice de facto de la mayoría de motores vectoriales modernos, incluido Qdrant. \citet{Johnson2019}, con FAISS, abordan el mismo problema desde una perspectiva distinta —el uso de GPUs para paralelizar tanto la búsqueda exhaustiva como la búsqueda aproximada mediante cuantización de producto—, demostrando que es posible construir un grafo de vecinos más cercanos sobre mil millones de vectores en menos de doce horas usando hardware especializado. \citet{Aumuller2020}, con ANN-Benchmarks, aportan un marco de evaluación empírica estandarizado que permite comparar de forma reproducible el desempeño de estos algoritmos (y de las bibliotecas que los implementan) bajo distintos conjuntos de datos y regímenes de precisión-velocidad, una metodología que este trabajo adapta en el Capítulo 3 para la comparación cuantitativa entre pgvector y Qdrant.

Estos algoritmos pueden implementarse embebidos en un motor de base de datos relacional de propósito general —extendido con un tipo de dato vectorial, como hace pgvector sobre PostgreSQL— o en un motor construido específicamente para datos vectoriales. \citet{Wang2021} describen con Milvus el diseño de esta segunda familia: una arquitectura que separa el almacenamiento del cómputo, soporta múltiples tipos de índice ANN de forma nativa y está optimizada para cargas de trabajo heterogéneas (CPU y GPU) a una escala que, según los autores, resulta difícil de sostener con extensiones vectoriales sobre motores relacionales no diseñados originalmente para ese patrón de acceso. La encuesta de \citet{Pan2024} sistematiza más de veinte sistemas de este tipo (incluyendo Milvus y Qdrant), proponiendo una taxonomía según su modelo de indexación, su soporte de búsqueda híbrida (vectorial más filtros estructurados) y su estrategia de escalamiento horizontal; mientras que la revisión de \citet{Taipalus2024} ofrece una introducción conceptual a los fundamentos, casos de uso y retos actuales de estos sistemas, señalando la consistencia transaccional y la actualización incremental del índice como puntos todavía activos de investigación.

El trabajo más directamente relevante para justificar una decisión de migración es el de \citet{Sanca2024}, quienes estudian empíricamente las rutas de acceso a datos en escenarios de búsqueda mixta vector-relacional —es decir, consultas que combinan similitud vectorial con filtros sobre atributos estructurados, exactamente el patrón de acceso de una plataforma de empleo que necesita combinar "semánticamente similar a la consulta" con "salario mayor a X" o "modalidad remota"—. Los autores muestran que, a medida que crece el volumen de datos y la selectividad de los filtros estructurados varía, las rutas de acceso disponibles en motores relacionales de propósito general presentan compromisos de rendimiento que un motor construido específicamente en torno a índices ANN puede evitar por diseño, al tratar la búsqueda vectorial como ciudadano de primera clase en lugar de como una extensión añadida sobre un modelo de datos relacional preexistente.

\textbf{Relación con JobAI.} JobAI almacenó inicialmente sus \textit{embeddings} —de 768 dimensiones, generados para cada fragmento de oferta laboral y cada currículo— directamente en PostgreSQL mediante pgvector, sin un índice HNSW o IVFFlat configurado explícitamente, lo que en la práctica implica una búsqueda por barrido secuencial: viable a decenas de miles de fragmentos, pero un compromiso real de escalabilidad ante un catálogo que crece cada seis horas con cada ciclo de ingesta. La evidencia revisada en esta sección —particularmente el hallazgo de \citet{Sanca2024} sobre búsqueda mixta vector-relacional, y la arquitectura de referencia de \citet{Wang2021}— fundamenta la decisión, documentada en el Capítulo 3, de migrar el almacenamiento vectorial a Qdrant como iteración de diseño dentro del ciclo de Design Science Research, y de conservar temporalmente el backend de pgvector para ejecutar la comparación cuantitativa de latencia y \textit{recall} entre ambos que se reporta en el Capítulo 5, empleando la metodología de benchmarking de \citet{Aumuller2020} como referencia.

\section{Design Science Research como paradigma metodológico}

Design Science Research (DSR) se distingue de la investigación puramente conductual en sistemas de información en que su objetivo no es explicar o predecir un fenómeno organizacional, sino diseñar y evaluar rigurosamente un artefacto que resuelve un problema relevante. \citet{Hevner2004} formalizan este paradigma mediante siete lineamientos —entre ellos, que el artefacto de diseño debe ser una contribución identificable (un constructo, modelo, método o instanciación), que su utilidad y calidad deben demostrarse mediante métodos de evaluación rigurosos, y que la investigación debe comunicarse tanto a audiencias orientadas a la tecnología como a la gestión—, y proponen un marco conceptual de dos ciclos complementarios: el ciclo de \textit{relevancia}, que conecta el artefacto con las necesidades del entorno de aplicación, y el ciclo de \textit{rigor}, que lo conecta con la base de conocimiento científico existente.

\citet{Peffers2007} operacionalizan estos principios en una metodología de seis actividades secuenciales pero no estrictamente lineales: identificación del problema y motivación, definición de los objetivos de una solución, diseño y desarrollo del artefacto, demostración de su uso en un contexto adecuado, evaluación de qué tan bien el artefacto resuelve el problema, y comunicación de los resultados. Un aporte relevante de esta metodología para el presente trabajo es que reconoce explícitamente distintos puntos de entrada al ciclo —incluyendo un punto de entrada centrado en el artefacto (\textit{design and development centered}), apropiado cuando la investigación parte de un artefacto ya existente y busca evaluarlo o mejorarlo, en lugar de partir de un problema completamente nuevo—.

El volumen editado por \citet{HevnerChatterjee2010} consolida, cinco años después del artículo seminal de 2004, tanto la teoría como casos de aplicación práctica de DSR, reafirmando el marco de los tres ciclos (relevancia, diseño y rigor) como la forma estándar de comunicar cómo un proyecto de DSR se nutre de la literatura existente (rigor) y retorna a ella una contribución evaluada (relevancia), a la vez que documenta el ciclo interno de construcción y evaluación iterativa del artefacto mismo.

\textbf{Relación con JobAI.} Este trabajo adopta el punto de entrada centrado en el artefacto de \citet{Peffers2007}: JobAI ya existe como instanciación funcional (el artefacto, en el sentido de \citealp{Hevner2004}), construida originalmente sin un proceso de evaluación académica formal. Las actividades de diseño y desarrollo de este ciclo de DSR corresponden a la migración del almacenamiento vectorial a Qdrant (Sección~\ref{sec:estado-arte-vectordb}); las actividades de demostración y evaluación corresponden a la ejecución del experimento comparativo contra un \textit{baseline} de palabra clave y a la medición de las métricas de fidelidad y precisión de recuperación (Capítulo 3); y la comunicación de resultados corresponde al presente documento. El Capítulo 3 desarrolla con mayor detalle esta correspondencia entre las actividades del modelo de \citet{Peffers2007} y las fases del cronograma de este proyecto.

\section{Clasificación híbrida de consultas y text-to-SQL}
\label{sec:estado-arte-hibrido}

No todas las preguntas que un usuario formula sobre un catálogo de ofertas laborales son preguntas de recuperación semántica. Una pregunta como \textit{"¿cuántas ofertas remotas de Python hay en la base de datos?"} tiene una respuesta exacta y verificable mediante una agregación SQL, y forzarla por una ruta de recuperación vectorial —que por diseño devuelve los fragmentos más similares, no un conteo exacto— produciría una respuesta aproximada o inconsistente. \citet{Biswal2024} formalizan esta observación con TAG (\textit{Table-Augmented Generation}), un paradigma que argumenta que tanto el text-to-SQL puro como el RAG puro cubren, cada uno por separado, solo un subconjunto estrecho de las preguntas que un usuario real querría formular sobre una base de datos: el primero se limita a lo expresable en álgebra relacional exacta, y el segundo se limita a preguntas resolubles mediante la recuperación de uno o pocos registros puntuales. Los autores proponen unificar ambos caminos —consulta exacta sobre datos estructurados y síntesis semántica sobre datos no estructurados— bajo un mismo marco de razonamiento.

La dificultad práctica de resolver la mitad "exacta" de ese problema mediante un LLM sin restricciones queda documentada por \citet{Li2023} con BIRD, un banco de pruebas de text-to-SQL a gran escala (12.751 pares de pregunta-consulta sobre 95 bases de datos reales de 37 dominios profesionales). Los autores muestran que incluso modelos de lenguaje muy capaces obtienen un desempeño sustancialmente inferior al humano al tener que generar SQL ejecutable y correcto sobre esquemas de bases de datos reales —con sus inconsistencias, ambigüedades y volúmenes de datos genuinos—, en contraste con bancos de pruebas anteriores que operaban sobre esquemas simplificados con pocos registros. Esta brecha entre el text-to-SQL de laboratorio y el text-to-SQL sobre datos reales motiva preferir, en un sistema de producción, una clasificación de la intención de la consulta antes de decidir si se genera SQL, en lugar de delegar esa decisión íntegramente a un LLM de propósito general.

\textbf{Relación con JobAI.} El clasificador de consultas de JobAI —descrito en el Capítulo 4— resuelve este mismo problema mediante un diseño pragmático de dos etapas: un conjunto de patrones \textit{regex} captura las señales léxicas inequívocas de intención de agregación ("cuántos", "promedio", "más común") y enruta esas consultas directamente a una plantilla SQL, evitando por completo el riesgo de generación de SQL abierta que documenta \citet{Li2023}; solo cuando ningún patrón coincide, un clasificador basado en LLM decide entre recuperación, agregación o una ruta híbrida que combina ambas en un mismo prompt. Este diseño puede leerse como una instanciación deliberadamente conservadora del espíritu de TAG \citep{Biswal2024} —reconocer que recuperación semántica y consulta exacta son necesarias para cubrir el espacio completo de preguntas de un usuario—, que privilegia la determinación y la seguridad de un conjunto acotado de plantillas SQL sobre la generalidad, más frágil, de un LLM generando SQL libremente sobre el esquema completo de la base de datos.

\section{Síntesis del marco de referencia}

Las seis líneas revisadas en este capítulo convergen en la premisa que motiva este trabajo: la limitación del emparejamiento léxico en la búsqueda de empleo (Secciones~\ref{sec:estado-arte-keyword} y siguientes) tiene, en la literatura de recuperación de información y de RAG, una solución arquitectónica madura y ampliamente adoptada por la industria del reclutamiento a gran escala; esa solución requiere, a su vez, decisiones de infraestructura (almacenamiento vectorial) y de enrutamiento de consultas (clasificación híbrida) que la literatura reciente permite fundamentar y no solo intuir. Lo que la literatura revisada no ofrece todavía —y que constituye el aporte específico de este trabajo— es una evaluación cuantitativa integrada de estas piezas sobre un artefacto real y en producción, siguiendo un proceso metodológico explícito de Design Science Research. El Capítulo 3 traduce esta síntesis en un diseño metodológico concreto.

    
    % Capitulo 3
    \chapter{Metodolog\'ia}
    \label{cap:metodologia}
    Este capítulo describe cómo se evalúa JobAI. Primero se enmarca el trabajo en el paradigma de Design Science Research y se reconstruyen los ciclos de diseño-evaluación por los que ya pasó el artefacto. Después se formulan las preguntas de evaluación y se detallan el sistema bajo prueba, el corpus, el conjunto de consultas, los sistemas comparados, el procedimiento de juicios de relevancia, las métricas, el diseño del benchmark de almacenes vectoriales y el análisis estadístico. El capítulo cierra con las amenazas a la validez y las medidas de reproducibilidad.

\section{Enfoque metodológico: Design Science Research}
\label{sec:met-dsr}

Este trabajo sigue el paradigma de Design Science Research (DSR) en la formulación de \citet{Hevner2004}, operacionalizada mediante el proceso de seis actividades de \citet{Peffers2007}. El artefacto es una instanciación de software, JobAI, que ya existía antes de iniciar esta investigación. Por ello se adopta el punto de entrada \emph{centrado en el diseño y desarrollo} que \citet{Peffers2007} prevén para la investigación que parte de un artefacto existente y busca evaluarlo y mejorarlo. La Tabla~\ref{tab:met-peffers} muestra la correspondencia entre las actividades del proceso y el contenido de este documento.

\begin{table}[htbp!]
    \centering
    \caption{Correspondencia entre las actividades de DSR \citep{Peffers2007} y este trabajo}
    \label{tab:met-peffers}
    \begin{tabular}{|p{4.2cm}|p{10cm}|}
        \hline
        \textbf{Actividad} & \textbf{Correspondencia en este trabajo} \\
        \hline
        Identificación del problema y motivación & Limitaciones del emparejamiento léxico en la búsqueda de empleo y ausencia de una evaluación cuantitativa de JobAI (Capítulo~\ref{cap:problema}). \\
        \hline
        Definición de los objetivos de la solución & Objetivo general y objetivos específicos (Capítulo~\ref{cap:problema}), traducidos en preguntas de evaluación medibles (Sección~\ref{sec:met-preguntas}). \\
        \hline
        Diseño y desarrollo & Arquitectura del pipeline RAG y su iteración actual: corrección de la consulta vectorial, migración a Qdrant y corrección de defectos detectados (Capítulo~\ref{cap:artefacto}). \\
        \hline
        Demostración & Ejecución del pipeline sobre un corpus real congelado y un conjunto de consultas de prueba (Secciones~\ref{sec:met-corpus} y~\ref{sec:met-consultas}). \\
        \hline
        Evaluación & Comparación contra búsqueda léxica, medición de fidelidad y latencia, y benchmark pgvector frente a Qdrant (Capítulo~\ref{cap:resultados}). \\
        \hline
        Comunicación & El presente documento. \\
        \hline
    \end{tabular}
\end{table}

\subsection{Ciclos de diseño-evaluación previos}
\label{sec:met-ciclos-previos}

El historial de versiones de JobAI documenta tres iteraciones de diseño-evaluación anteriores a este trabajo, realizadas en abril de 2026 sobre un corpus distinto al actual: unas 3.200 ofertas, mayoritariamente del agregador Adzuna, que dejó de usarse después. Cada iteración se evaluó con 15 consultas y un modelo de lenguaje como juez de relevancia. La Tabla~\ref{tab:met-iteraciones} resume los resultados registrados.

\begin{table}[htbp!]
    \centering
    \caption{Iteraciones previas de diseño-evaluación registradas en el historial de JobAI (15 consultas, $k=5$)}
    \label{tab:met-iteraciones}
    \begin{tabular}{|p{6.2cm}|p{2.6cm}|c|c|c|}
        \hline
        \textbf{Configuración} & \textbf{Juez} & \textbf{P@5} & \textbf{MRR} & \textbf{nDCG} \\
        \hline
        Fragmentación por ventana deslizante & Llama 3.1 8B & 66,7\,\% & 0,656 & 0,702 \\
        \hline
        Fragmentación estructurada (4 tipos) & Llama 3.1 8B & 92,0\,\% & 0,967 & 0,967 \\
        \hline
        Estructurada + ruta de agregación & Llama 3.1 8B & 90,7\,\% & 1,000 & 0,995 \\
        \hline
        Estructurada + ruta de agregación & Modelo comercial & 98,7\,\% & 1,000 & 0,994 \\
        \hline
    \end{tabular}
\end{table}

Estas iteraciones justificaron decisiones de diseño concretas, como la fragmentación estructurada que el Capítulo~\ref{cap:artefacto} describe. Su valor como evidencia, sin embargo, es limitado por cuatro razones:
\begin{itemize}
    \item El conjunto de consultas es pequeño y se diseñó a la vista del corpus.
    \item No existe un sistema de referencia contra el cual comparar.
    \item El juez de relevancia nunca se validó contra juicios humanos.
    \item El corpus ya no corresponde al que la plataforma sirve hoy.
\end{itemize}
La última fila de la Tabla~\ref{tab:met-iteraciones} ilustra la tercera limitación: la misma configuración obtiene una precisión de 90,7\,\% o de 98,7\,\% según el modelo que actúe como juez. La evaluación diseñada en este capítulo responde directamente a estas cuatro limitaciones y constituye, en términos de DSR, la cuarta iteración del ciclo.

\section{Preguntas de evaluación}
\label{sec:met-preguntas}

Los objetivos específicos tercero y cuarto se operacionalizan en cuatro preguntas de evaluación (PE):

\begin{description}
    \item[PE1 (calidad de recuperación).] ¿La recuperación densa de JobAI obtiene mayor precisión y mejor ordenamiento de resultados que la búsqueda léxica, y en qué tipos de consulta se concentra la diferencia?
    \item[PE2 (fidelidad de la generación).] ¿En qué medida las respuestas que genera el pipeline se sostienen en las ofertas recuperadas y responden a la pregunta formulada?
    \item[PE3 (latencia).] ¿Cuánto tarda cada etapa del pipeline (\textit{embedding} de la consulta, recuperación, construcción del contexto y generación) y cuál es la latencia percibida por el usuario?
    \item[PE4 (almacenamiento vectorial).] ¿Cómo se comparan pgvector y Qdrant en la relación entre \textit{recall} y latencia, con y sin filtros estructurados y a distintas escalas de datos?
\end{description}

\section{Sistema bajo evaluación}
\label{sec:met-sistema}

Se evalúa una configuración congelada de JobAI, identificada por el \textit{commit} del repositorio que acompaña a cada corrida de evaluación. La Tabla~\ref{tab:met-config} resume sus parámetros.

\begin{table}[htbp!]
    \centering
    \caption{Configuración del sistema bajo evaluación}
    \label{tab:met-config}
    \begin{tabular}{|p{5cm}|p{9.2cm}|}
        \hline
        \textbf{Componente} & \textbf{Configuración} \\
        \hline
        Modelo de \textit{embeddings} & e5-base-v2 \citep{Wang2022E5}, 768 dimensiones, ejecutado localmente; prefijos \texttt{query:} y \texttt{passage:} \\
        \hline
        Fragmentación & Estructurada (identidad, requisitos, responsabilidades, beneficios y salario); ventana deslizante de 800 caracteres con solapamiento de 200 cuando el análisis estructurado falla \\
        \hline
        Recuperación & 15 fragmentos más similares por coseno, agrupados por oferta; se conservan las 5 ofertas con mayor similitud \\
        \hline
        Índice vectorial & HNSW con $m=16$ y $ef_{construction}=128$ en pgvector y en Qdrant; $ef_{search}=40$ \\
        \hline
        Modelo generador & \texttt{qwen/qwen3.8-27b} servido por Groq, temperatura 0 durante la evaluación \\
        \hline
        Clasificación de consultas & Patrones \textit{regex} más un clasificador basado en LLM restringido a un catálogo cerrado de plantillas SQL \\
        \hline
    \end{tabular}
\end{table}

En producción la temperatura de generación es 0,2. Para la evaluación se fija en 0 y todas las respuestas se almacenan en caché, de modo que una nueva ejecución del análisis no exige regenerarlas. Las consultas se envían sin usuario autenticado; de lo contrario, la personalización por currículo alteraría la consulta y confundiría el efecto de la recuperación con el del perfil del usuario.

\section{Corpus de evaluación}
\label{sec:met-corpus}

El corpus se obtuvo el 10 de octubre de 2026 mediante una ejecución completa del proceso de ingesta de JobAI sobre sus tres fuentes públicas: Remotive, Jobicy y Findwork. El proceso se ejecutó sobre una base de datos dedicada y con el programador de ingesta periódica desactivado, para que el catálogo no cambiara durante la evaluación. El resultado se congeló en una copia de seguridad (\textit{snapshot}) identificada por su resumen SHA-256. Todas las corridas de evaluación parten de esa copia.

El corpus contiene 1.068 ofertas: 895 de Findwork, 162 de Jobicy y 11 de Remotive. De ellas, 682 están marcadas como remotas y 504 informan salario. Las ofertas se dividen en 4.013 fragmentos, todos con su \textit{embedding}, con una longitud media de 467 caracteres. El análisis estructurado con el modelo local Llama 3.1 8B falló por completo en 25 ofertas (2,3\,\%), a las que se aplicó la fragmentación por ventana deslizante, y fue parcial, sin resumen, en otras 231.

El corpus tiene tres sesgos conocidos, que se tienen en cuenta al interpretar los resultados:
\begin{itemize}
    \item Los conectores de Remotive y Jobicy descartan las ofertas sin salario explícito, lo que explica su baja participación.
    \item La API de Findwork interrumpió la descarga al alcanzar su límite de solicitudes, y aun así aporta el 84\,\% del corpus, de modo que este está dominado por una sola fuente.
    \item El normalizador de salarios del artefacto no convierte monedas ni algunos montos por hora, por lo que unas pocas ofertas tienen salarios atípicos. La mediana de los salarios informados es de 175.000 dólares anuales.
\end{itemize}

\section{Conjunto de consultas de prueba}
\label{sec:met-consultas}

El conjunto de prueba contiene 75 consultas en inglés, que es el idioma de las ofertas, distribuidas en cinco categorías de 15 consultas cada una:

\begin{description}
    \item[Exactas.] Consultas cuyos términos aparecen literalmente en las ofertas relevantes, por ejemplo \textit{"senior react developer"}. Es el caso favorable a la búsqueda léxica.
    \item[Desajuste de vocabulario.] Consultas que describen la necesidad con términos distintos a los de las ofertas: sinónimos, títulos equivalentes o descripciones funcionales. Es el caso que motiva este trabajo (Sección~\ref{sec:problema}).
    \item[Con restricciones.] Consultas que combinan un perfil con condiciones sobre atributos estructurados, como salario, modalidad o tipo de contrato.
    \item[Ruidosas.] Consultas con errores ortográficos, abreviaturas o formulación vaga.
    \item[De agregación.] Preguntas estadísticas sobre el catálogo, por ejemplo cuántas ofertas remotas de Python existen. Se evalúan aparte (Sección~\ref{sec:met-metricas}).
\end{description}

Cada consulta se asocia a una \emph{necesidad de información}, es decir, una descripción breve de lo que el usuario busca. A partir de la distribución real de títulos y habilidades del corpus se redactan además 25 consultas en español, traducciones paralelas de consultas en inglés que comparten su necesidad de información. Como la relevancia de una oferta depende de la necesidad y no del idioma en que se formula, las consultas en español reutilizan los juicios de relevancia de su contraparte en inglés. Este subconjunto mide la recuperación entre idiomas, con una salvedad que se anticipa: e5-base-v2 es un modelo entrenado en inglés, por lo que se espera un desempeño inferior en español.

\section{Sistemas comparados}
\label{sec:met-sistemas}

Se comparan cinco sistemas de recuperación sobre el mismo corpus. Todos producen una lista ordenada de 10 ofertas por consulta:

\begin{description}
    \item[ILIKE.] La búsqueda por palabra clave que la plataforma ofrece hoy en su panel de filtros: coincidencia de subcadena, sin distinguir mayúsculas, sobre el título y la empresa. Representa la línea base \emph{de la propia plataforma}.
    \item[BM25.] El modelo probabilístico de relevancia Okapi BM25 \citep{Robertson2009}, con $k_1 = 1{,}2$ y $b = 0{,}75$, sobre el título, la descripción y las habilidades de cada oferta. Representa la línea base léxica \emph{fuerte} de la literatura de recuperación de información. Su inclusión evita que la comparación se haga únicamente contra un sistema débil.
    \item[Denso, pgvector exacto.] Búsqueda exhaustiva por similitud coseno en PostgreSQL. Sirve de referencia para medir el efecto de la búsqueda aproximada.
    \item[Denso, pgvector HNSW.] La misma búsqueda servida por el índice HNSW de pgvector.
    \item[Denso, Qdrant HNSW.] La misma búsqueda servida por Qdrant con parámetros de índice idénticos.
\end{description}

En los tres sistemas densos se recuperan los 50 fragmentos más similares a la consulta y se agregan por oferta, tomando como puntuación de cada oferta la máxima similitud entre sus fragmentos. Es la misma regla de agregación que usa el pipeline en producción. Así la unidad evaluada es la oferta, que es lo que el usuario ve, y no el fragmento.

\section{Juicios de relevancia}
\label{sec:met-juicios}

\subsection{Agrupamiento de resultados}

Los juicios se construyen mediante agrupamiento (\textit{pooling}) de profundidad 10, el procedimiento estándar de las colecciones de prueba en recuperación de información \citep{Sanderson2010}. Para cada consulta se juzga la unión de las 10 primeras ofertas devueltas por cada sistema. Las ofertas que ningún sistema recupera en esas posiciones se consideran no relevantes, una aproximación cuyo efecto se discute en la Sección~\ref{sec:met-validez}.

\subsection{Modelo de lenguaje como juez}

Cada par consulta-oferta del agrupamiento recibe una etiqueta graduada en la escala de cuatro niveles de las evaluaciones TREC, asignada por un modelo de lenguaje:
\begin{itemize}
    \item 0: no relevante.
    \item 1: relacionada.
    \item 2: altamente relevante.
    \item 3: perfectamente relevante.
\end{itemize}
El prompt sigue el diseño de UMBRELA \citep{Upadhyay2024}, una reproducción abierta del procedimiento con el que \citet{Thomas2024} mostraron que un LLM puede predecir las preferencias de evaluadores humanos con una precisión comparable a la de estos. El juez recibe la necesidad de información y una ficha estructurada de la oferta: título, empresa, ubicación, modalidad, salario, habilidades, resumen y un extracto de la descripción.

El uso de un LLM como juez tiene riesgos documentados, entre ellos sesgos sistemáticos, sensibilidad al prompt y posible circularidad cuando el juez y el sistema evaluado comparten modelo \citep{Faggioli2023,Zheng2023}. Para mitigarlos se toman tres medidas:
\begin{itemize}
    \item El juez es un modelo distinto del generador de JobAI.
    \item Durante el desarrollo del arnés de evaluación se usa un modelo local (Llama 3.1 8B), y los resultados definitivos se obtienen con un modelo comercial de mayor capacidad (\pendiente{modelo del juez final}).
    \item Las etiquetas del juez se validan contra juicios humanos, como se describe a continuación.
\end{itemize}

\subsection{Validación humana del juez}

Del conjunto de pares juzgados se extrae una muestra aleatoria del 20\,\%, estratificada por categoría de consulta y por etiqueta del juez, con semilla fija. El autor la etiqueta a ciegas, sin conocer la etiqueta del modelo. La muestra se usa en dos partes:
\begin{itemize}
    \item Una cuarta parte (5\,\% del total) se usa para calibrar el prompt del juez.
    \item Las tres cuartas partes restantes (15\,\% del total) se reservan para medir el acuerdo final, sin volver a modificar el prompt.
\end{itemize}

El acuerdo se mide con el coeficiente kappa de \citet{Cohen1960}, en dos variantes: sobre la relevancia binaria (etiqueta $\geq 2$) y ponderado cuadráticamente sobre la escala completa. Se reporta con un intervalo de confianza bootstrap. Siguiendo la escala de interpretación de \citet{Landis1977}, se considera aceptable un acuerdo al menos sustancial ($\kappa > 0{,}60$).

\section{Métricas}
\label{sec:met-metricas}

\subsection{Calidad de recuperación (PE1)}

Sobre las listas de 10 ofertas se calculan cinco métricas:
\begin{itemize}
    \item La precisión en las posiciones 5 y 10 (P@5, P@10).
    \item La ganancia acumulada descontada normalizada en la posición 10 (nDCG@10) con las etiquetas graduadas \citep{Jarvelin2002}.
    \item El rango recíproco medio (MRR).
    \item El \textit{recall} en la posición 10, relativo al conjunto de ofertas relevantes identificadas en el agrupamiento.
\end{itemize}
Las métricas binarias consideran relevante una oferta con etiqueta $\geq 2$. El cálculo se realiza con la biblioteca \texttt{ir\_measures} \citep{MacAvaney2022}, que unifica las implementaciones de referencia de la comunidad. Los resultados se desagregan por categoría de consulta y por idioma.

Las consultas de agregación no producen una lista de ofertas. Para ellas se miden dos cosas: la exactitud del enrutamiento, es decir, si el clasificador elige la ruta y la plantilla SQL esperadas, y la exactitud de la respuesta frente al resultado de una consulta SQL de referencia escrita a mano.

\subsection{Fidelidad de la generación (PE2)}

Para las consultas de recuperación se generan respuestas mediante la ruta real del pipeline y se calculan dos métricas de RAGAS \citep{Es2024}, empleando la implementación de referencia de los autores:

\begin{itemize}
    \item \textbf{Fidelidad} (\textit{faithfulness}). Un LLM descompone la respuesta en afirmaciones atómicas y verifica cada una contra el contexto recuperado:
    \[
    F = \frac{|V|}{|S|},
    \]
    donde $S$ es el conjunto de afirmaciones de la respuesta y $V \subseteq S$ el de las afirmaciones sostenidas por el contexto.
    \item \textbf{Relevancia de la respuesta} (\textit{answer relevancy}). Un LLM genera $n$ preguntas a partir de la respuesta, y la métrica es la similitud coseno media entre sus \textit{embeddings} y el de la pregunta original:
    \[
    AR = \frac{1}{n}\sum_{i=1}^{n} \cos(\mathbf{e}_q, \mathbf{e}_{q_i}).
    \]
\end{itemize}

\subsection{Latencia (PE3)}

El pipeline se instrumenta para registrar la duración de cada etapa: \textit{embedding} de la consulta, búsqueda vectorial, lectura de las ofertas y construcción del contexto, tiempo hasta el primer \textit{token} generado (TTFT) y tiempo total. Las etapas locales se miden con 30 repeticiones por consulta y por almacén vectorial, tras 5 de calentamiento. La generación, que depende de un servicio externo, se mide con 3 repeticiones sobre un subconjunto de 30 consultas. Se reportan la mediana (p50) y el percentil 95 (p95). La latencia del servicio externo se reporta por separado de la de los componentes propios del artefacto.

\section{Benchmark de almacenes vectoriales (PE4)}
\label{sec:met-benchmark}

La comparación entre pgvector y Qdrant adapta la metodología de ANN-Benchmarks \citep{Aumuller2020}: para cada configuración del parámetro de búsqueda $ef_{search} \in \{16, 32, 64, 128, 256, 512\}$ se mide el \textit{recall}@10 frente a la búsqueda exacta y la latencia de las consultas. El resultado se presenta como curvas de compromiso entre \textit{recall} y latencia, en lugar de un único punto de operación.

\subsection{Datos}

El corpus real es pequeño: unos pocos miles de fragmentos. A esa escala es esperable que ambos motores rindan de forma casi idéntica. Para estudiar el comportamiento a mayor escala se generan conjuntos sintéticos de 100.000 y 500.000 vectores, y de un millón si los recursos de cómputo lo permiten. Cada vector sintético se obtiene interpolando dos vectores reales del corpus y agregando ruido gaussiano escalado por la desviación estándar de cada dimensión, para luego normalizarlo. La magnitud del ruido se calibra de modo que la distribución de similitudes al vecino más cercano se parezca a la del corpus real; la diferencia se cuantifica con el estadístico de Kolmogórov-Smirnov. Cada vector sintético hereda los atributos estructurados de la oferta de la que proviene, de modo que los filtros conservan su correlación con el contenido.

\subsection{Protocolo}

Se emplean unas 1.000 consultas: las del conjunto de prueba más títulos de ofertas del corpus. La verdad de referencia son los 100 vecinos más cercanos exactos, calculados por fuerza bruta, también bajo cada filtro. Para reproducir el patrón de búsqueda mixta vector-relacional estudiado por \citet{Sanca2024}, cada configuración se mide sin filtro y con filtros estructurados de selectividad 1\,\%, 10\,\% y 50\,\%. En pgvector se comparan además los barridos iterativos del índice (\texttt{hnsw.\allowbreak iterative\_scan}), disponibles desde la versión 0.8, frente a su ausencia.

Ambos motores usan el mismo índice HNSW ($m=16$, $ef_{construction}=128$) y se ejecutan uno a la vez, en contenedores con los mismos límites de CPU y memoria. Para cada configuración se registran siete medidas:
\begin{itemize}
    \item \textit{recall}@10;
    \item latencia p50, p95 y p99;
    \item consultas por segundo con 1 y con 8 clientes concurrentes;
    \item completitud, es decir, si se devolvieron $k$ resultados;
    \item tiempo de construcción del índice;
    \item tamaño en disco;
    \item memoria pico.
\end{itemize}
Cada medición se repite tres veces y se reporta la mediana.

\section{Análisis estadístico}
\label{sec:met-estadistica}

Las comparaciones entre sistemas de recuperación son pareadas: todos los sistemas responden las mismas consultas. Siguiendo la comparación de pruebas de significancia en recuperación de información de \citet{Smucker2007}, que desaconseja la prueba de signos y la de Wilcoxon en este contexto, la diferencia en cada métrica se contrasta con la prueba de aleatorización pareada, con 10.000 permutaciones. Se acompaña de un intervalo de confianza bootstrap del 95\,\% para la diferencia de medias. Como se realizan varias comparaciones simultáneas (cada sistema denso frente a cada línea base, por métrica), los valores $p$ se ajustan con el procedimiento secuencial de \citet{Holm1979}. El nivel de significancia es $\alpha = 0{,}05$.

\section{Amenazas a la validez}
\label{sec:met-validez}

\begin{description}
    \item[Validez interna.] Hay tres amenazas principales.
    \begin{itemize}
        \item \emph{El juez automático puede tener sesgos.} Se mitiga con la validación humana y el reporte del kappa.
        \item \emph{El agrupamiento puede favorecer a los sistemas que lo originaron.} Si se añade un sistema después, solo se juzgan sus resultados no evaluados previamente.
        \item \emph{El autor redacta las consultas.} Se mitiga construyéndolas a partir de la distribución real del corpus y por categorías definidas de antemano.
    \end{itemize}
    \item[Validez externa.] Los resultados corresponden a un dominio (empleo remoto, mayoritariamente tecnológico), a un idioma de las ofertas (inglés) y a un único modelo de \textit{embeddings}. No se generalizan a otros dominios ni modelos. El tamaño del corpus real limita las conclusiones sobre escalabilidad; el benchmark sintético las complementa, pero no las sustituye.
    \item[Medición de latencia.] La latencia se mide desde el cliente, de modo que incluye el protocolo de comunicación y la biblioteca cliente de cada motor (\textit{wire protocol} de PostgreSQL frente a gRPC en Qdrant). Es la latencia que percibe una aplicación, pero a escalas pequeñas puede estar dominada por esa sobrecarga y no por el motor. Por eso los resultados se interpretan junto con la escala del conjunto de datos.
    \item[Validez de constructo.] La relevancia juzgada, ya sea por un modelo o por el autor, aproxima pero no equivale a la utilidad que un usuario real percibiría. La evaluación con usuarios queda fuera del alcance (Sección~\ref{sec:problema}).
    \item[Confiabilidad.] Los modelos de lenguaje no son deterministas, incluso a temperatura 0. Además, los proveedores retiran modelos sin previo aviso: durante este trabajo dejaron de estar disponibles los dos modelos generadores configurados originalmente en JobAI (\texttt{qwen3-32b} y \texttt{qwen3.6-27b}), y se adoptó su sucesor, \texttt{qwen3.8-27b}. Ambos riesgos se mitigan registrando cada respuesta y cada veredicto con el identificador exacto del modelo y la fecha, y ejecutando el análisis sobre esa caché.
\end{description}

\section{Reproducibilidad}
\label{sec:met-reproducibilidad}

Cada corrida de evaluación se ejecuta mediante una interfaz de línea de comandos del propio artefacto y escribe sus resultados en un directorio propio con cuatro elementos:
\begin{itemize}
    \item el \textit{commit} del código;
    \item los modelos y parámetros empleados;
    \item el resumen SHA-256 del corpus;
    \item los resultados por consulta.
\end{itemize}
Las versiones de PostgreSQL y pgvector (0.8.2), de Qdrant (1.19.1) y de las bibliotecas de software están fijadas. Los veredictos del juez se almacenan en una caché indexada por modelo, versión del prompt, necesidad de información y oferta, de modo que el análisis puede repetirse sin nuevas llamadas a modelos de lenguaje. El código, el conjunto de consultas, los juicios de relevancia y los resultados se publican junto con este documento. El corpus no se publica, para respetar los términos de uso de las fuentes, pero el procedimiento para reconstruirlo está documentado.

    
    % Capitulo 4
    \chapter{Dise\~no y desarrollo del artefacto}
    \label{cap:artefacto}
    Este capítulo documenta el artefacto evaluado. Las primeras secciones describen la arquitectura de JobAI, el proceso de ingesta y el pipeline de consulta en su configuración de referencia, que es la que se evalúa en el Capítulo~\ref{cap:resultados}. Las secciones siguientes corresponden a las dos iteraciones de diseño de este trabajo. La primera, de \emph{corrección}, hizo que el artefacto implementara lo que su documentación describe y puso el almacenamiento vectorial detrás de una interfaz común. La segunda, de \emph{mejora}, incorporó a la recuperación filtros estructurados, búsqueda híbrida, reordenamiento y un modelo de \textit{embeddings} multilingüe. La última sección reúne las fallas que la evaluación reveló en ambas iteraciones y las clasifica según su efecto.

\section{Visión general de la arquitectura}
\label{sec:art-arquitectura}

JobAI es una plataforma de agregación de ofertas de empleo remoto cuyo componente central es un copiloto conversacional basado en RAG. El sistema se organiza en tres capas (Figura~\ref{fig:art-arquitectura}):
\begin{itemize}
    \item un cliente web en Next.js, responsable de la autenticación, la navegación del catálogo y el panel conversacional;
    \item un servidor en NestJS, que orquesta la ingesta, la recuperación y la generación;
    \item una capa de datos compuesta por PostgreSQL, como almacén relacional, y un almacén vectorial que, tras la iteración descrita en la Sección~\ref{sec:art-vectorial}, puede ser pgvector o Qdrant.
\end{itemize}
Los modelos se consumen como servicios. Un servidor propio, ejecutado localmente sobre una GPU, sirve el modelo de \textit{embeddings} y el modelo de reordenamiento; los modelos de lenguaje se consumen a través de un proveedor externo, Groq.

Este capítulo describe la \emph{configuración de referencia} del artefacto, definida en el repositorio junto con el código. El despliegue público de la plataforma conserva todavía la configuración anterior a la Iteración~2, con un proveedor externo de \textit{embeddings} y sin reordenamiento, y se actualizará a la configuración de referencia.

\begin{figure}[htbp!]
    \centering
    \caption{Arquitectura de JobAI}
    \label{fig:art-arquitectura}
    \begin{tikzpicture}[
        box/.style={draw, rounded corners, align=center, minimum height=1cm, text width=3.2cm, font=\small},
        store/.style={draw, cylinder, shape border rotate=90, aspect=0.15, align=center, minimum height=1.2cm, text width=2.4cm, font=\small},
        arr/.style={->, >=stealth}
    ]
        \node[box, text width=2.2cm] (fe) at (0,0) {Cliente web\\(Next.js)};
        \node[box, text width=2.2cm] (be) at (5.4,0) {Servidor\\(NestJS)};
        \node[box] (models) at (10.2,2.1) {Modelos de lenguaje\\(Groq)};
        \node[box] (local) at (10.2,0) {\textit{Embeddings} y reordenamiento (servidor local, GPU)};
        \node[box] (sources) at (10.2,-2.1) {Fuentes de ofertas\\(Remotive, Jobicy, Findwork)};
        \node[store] (pg) at (3.8,-3.2) {PostgreSQL\\(+ pgvector)};
        \node[store] (qd) at (7.0,-3.2) {Qdrant};
        \draw[arr] (fe) -- node[above, font=\scriptsize] {REST / NDJSON} (be);
        \draw[arr] (be) -- (models);
        \draw[arr] (be) -- (local);
        \draw[arr] (sources) -- (be);
        \draw[arr] (be) -- (pg);
        \draw[arr, dashed] (be) -- (qd);
    \end{tikzpicture}
\end{figure}

El servidor se organiza en módulos con responsabilidades separadas (Tabla~\ref{tab:art-modulos}). La consulta conversacional se expone mediante un \textit{endpoint} que transmite la respuesta de forma incremental como una secuencia de objetos JSON delimitados por saltos de línea (NDJSON). Los eventos son de cuatro tipos: inicio, \textit{token}, fin con las fuentes citadas y error.

\begin{table}[htbp!]
    \centering
    \caption{Módulos del servidor de JobAI}
    \label{tab:art-modulos}
    \begin{tabular}{|p{3.2cm}|p{11cm}|}
        \hline
        \textbf{Módulo} & \textbf{Responsabilidad} \\
        \hline
        Ingesta & Obtención periódica de ofertas, deduplicación, análisis estructurado, normalización de salarios y fragmentación \\
        \hline
        \textit{Embeddings} & Generación de vectores de consultas y fragmentos, con caché por modelo \\
        \hline
        Consulta & Clasificación de la consulta, extracción de restricciones y orquestación de las rutas de recuperación, agregación e híbrida \\
        \hline
        Recuperación & Búsqueda densa y léxica bajo el mismo filtro, fusión de rankings y reordenamiento \\
        \hline
        RAG & Construcción del contexto y del prompt \\
        \hline
        LLM & Acceso a los modelos generativos, con reintentos y rotación de claves \\
        \hline
        Almacenamiento & Repositorios relacionales y almacenes vectoriales (pgvector y Qdrant) \\
        \hline
        Currículo & Carga, extracción de texto y análisis estructurado del currículo del usuario \\
        \hline
        Autenticación y usuario & Tokens JWT con rotación, perfil y ofertas guardadas \\
        \hline
    \end{tabular}
\end{table}

\section{Ingesta de datos}
\label{sec:art-ingesta}

Un proceso programado consulta las tres fuentes públicas cada seis horas. Puede desactivarse mediante configuración, como se hizo para congelar el corpus de evaluación. Cada ciclo sigue cuatro pasos:
\begin{enumerate}
    \item \textbf{Deduplicación.} Cada oferta se identifica mediante el resumen SHA-256 de su título, empresa, descripción y ubicación. Las ofertas ya almacenadas y con \textit{embeddings} se omiten.
    \item \textbf{Análisis estructurado.} Las ofertas nuevas pasan por un modelo de lenguaje que extrae en formato JSON los siguientes campos: resumen, salario, responsabilidades, requisitos, beneficios, habilidades, tipo de contrato y modalidad. El modelo es configurable. La configuración de referencia usa \texttt{gpt-oss-120b} vía Groq, pero el corpus de evaluación se procesó con Llama 3.1 8B ejecutado localmente, de modo que la evaluación corresponde a fichas generadas por ese modelo (Sección~\ref{sec:met-corpus}). El prompt instruye explícitamente al modelo a devolver \texttt{null} cuando un campo no aparece en el texto, en lugar de inferirlo. Los salarios se normalizan a dólares anuales; los valores fuera del intervalo de 10.000 a 1.000.000 dólares, o con un mínimo mayor que el máximo, se registran como no informados, sin intentar adivinar la moneda o el periodo.
    \item \textbf{Fragmentación.} Si el análisis estructurado produce al menos un campo, la oferta se divide en fragmentos \emph{tipados}:
    \begin{itemize}
        \item un fragmento de identidad (título, empresa, ubicación, modalidad y resumen);
        \item uno de requisitos;
        \item uno de responsabilidades y habilidades;
        \item uno de beneficios y salario.
    \end{itemize}
    Si el análisis falla, se aplica una ventana deslizante de 800 caracteres con solapamiento de 200, que respeta los límites de párrafo, viñeta y oración, con un máximo de cuatro ventanas por oferta.
    \item \textbf{\textit{Embeddings} y almacenamiento.} Cada fragmento se convierte en un vector de 768 dimensiones con el modelo multilingual-e5-base \citep{Wang2024mE5}, con el prefijo \texttt{passage:} que el modelo espera para los documentos, y se almacena junto con su texto, su tipo y el modelo que lo produjo. El servidor rechaza vectores de un modelo distinto al configurado, de modo que una consulta nunca se compara con fragmentos codificados por otro modelo.
\end{enumerate}
Al terminar cada ciclo se reconstruye el índice léxico que usa la búsqueda híbrida (Sección~\ref{sec:art-recuperacion}).

La fragmentación tipada responde a un hallazgo de una iteración previa (Sección~\ref{sec:met-ciclos-previos}). Un fragmento que mezcla requisitos, beneficios y descripción compite peor en la recuperación que un fragmento que contiene solo la información por la que se pregunta.

\section{Pipeline de consulta}
\label{sec:art-consulta}

La Figura~\ref{fig:art-pipeline} resume el recorrido de una consulta.

\begin{figure}[htbp!]
    \centering
    \caption{Pipeline de consulta de JobAI}
    \label{fig:art-pipeline}
    \begin{tikzpicture}[
        stage/.style={draw, rounded corners, align=center, text width=4.2cm, minimum height=0.9cm, font=\small},
        half/.style={stage, text width=3.2cm},
        decision/.style={draw, diamond, aspect=2.2, align=center, inner sep=1pt, font=\small},
        arr/.style={->, >=stealth}
    ]
        \node[stage] (q) at (0,0) {Consulta del usuario};
        \node[decision] (cls) at (0,-1.6) {Clasificación};
        \node[stage] (rag) at (-3.6,-3.2) {Ruta de recuperación\\(con restricciones como filtro)};
        \node[stage] (agg) at (3.6,-3.2) {Ruta de agregación\\(plantilla SQL)};
        \node[half] (dense) at (-5.5,-4.8) {Búsqueda densa: 50 fragmentos $\rightarrow$ ofertas};
        \node[half] (bm25) at (-1.7,-4.8) {BM25:\\50 ofertas};
        \node[stage] (rrf) at (-3.6,-6.4) {Fusión RRF\\(20 candidatas)};
        \node[stage] (rr) at (-3.6,-7.8) {Reordenamiento\\(5 ofertas)};
        \node[stage] (ctx) at (-3.6,-9.2) {Contexto por oferta + campos estructurados};
        \node[stage] (sum) at (3.6,-4.8) {Resumen de la agregación};
        \node[stage] (llm) at (0,-10.8) {Prompt + generación (LLM)};
        \node[stage] (out) at (0,-12.2) {Respuesta incremental con fuentes};
        \draw[arr] (q) -- (cls);
        \draw[arr] (cls) -| node[above, pos=0.25, font=\scriptsize] {recuperación} (rag);
        \draw[arr] (cls) -| node[above, pos=0.25, font=\scriptsize] {agregación} (agg);
        \draw[arr] (rag) -- (dense);
        \draw[arr] (rag) -- (bm25);
        \draw[arr] (dense) -- (rrf);
        \draw[arr] (bm25) -- (rrf);
        \draw[arr] (rrf) -- (rr);
        \draw[arr] (rr) -- (ctx);
        \draw[arr] (agg) -- (sum);
        \draw[arr] (ctx) |- (llm);
        \draw[arr] (sum) |- (llm);
        \draw[arr] (llm) -- (out);
    \end{tikzpicture}
\end{figure}

\subsection{Clasificación de la consulta}
\label{sec:art-clasificacion}

Un conjunto de patrones \textit{regex} detecta señales léxicas de tres tipos:
\begin{itemize}
    \item de recuperación (\textit{find}, \textit{looking for}, \textit{jobs with});
    \item de agregación (\textit{how many}, \textit{average}, \textit{salary}, \textit{list all});
    \item de restricción, en inglés y en español: montos de salario, modalidad (\textit{remote}, \textit{remoto}, \textit{presencial}) y tipo de contrato (\textit{contract}, \textit{tiempo completo}, \textit{freelance}, \textit{prácticas}).
\end{itemize}
Las consultas con señales de recuperación y sin señales de agregación ni de restricción van directamente a la ruta de recuperación, sin invocar ningún modelo de lenguaje. En los demás casos, un clasificador basado en LLM, con temperatura 0, decide entre recuperación, agregación o una ruta híbrida. Según el caso, el clasificador selecciona además una de 24 plantillas SQL parametrizadas, como conteos, estadísticas salariales o listados por habilidad, o extrae las restricciones de la consulta: salario mínimo o máximo en dólares anuales, modalidad remota y tipo de contrato. El modelo nunca genera SQL libre, y las consultas de agregación se ejecutan con una conexión de solo lectura.

Sobre la respuesta del clasificador se aplican tres reglas de enrutamiento deterministas:
\begin{itemize}
    \item una consulta con restricciones extraídas se atiende por la ruta de recuperación, con las restricciones como filtro;
    \item las plantillas que responden con estadísticas, valores distintos u ofertas recientes se atienden siempre por la ruta de agregación, aunque el clasificador haya etiquetado la consulta como recuperación;
    \item una consulta de agregación sin una plantilla válida recurre a la ruta de recuperación.
\end{itemize}

El título del cargo y las ubicaciones distintas de \textit{remoto} no se convierten en filtros: exigir su coincidencia literal reproduciría la limitación léxica que el pipeline busca superar, así que se dejan a la recuperación semántica. Como las fuentes suelen etiquetar el trabajo independiente como contrato, una consulta por trabajo \textit{freelance} acepta ambos tipos. Los filtros que el usuario fija explícitamente en el panel de búsqueda tienen prioridad sobre los extraídos del texto.

\subsection{Recuperación}
\label{sec:art-recuperacion}

La recuperación combina dos búsquedas de primera etapa, una fusión de rankings y un reordenamiento:
\begin{enumerate}
    \item \textbf{Búsqueda densa.} La consulta se convierte en un vector con el prefijo \texttt{query:}. Si el usuario cargó un currículo, se antepone a la consulta un perfil breve con sus tres cargos más recientes y sus diez habilidades principales. Se buscan los 50 fragmentos más similares por coseno que cumplen el filtro; se agrupan por oferta y cada oferta toma la similitud de su mejor fragmento.
    \item \textbf{Búsqueda léxica.} En paralelo, BM25 \citep{Robertson2009}, con $k_1 = 1{,}2$ y $b = 0{,}75$, busca las 50 ofertas que mejor coinciden con el texto de la consulta entre las que cumplen el mismo filtro. El índice se mantiene en memoria sobre el título, las habilidades y la descripción de cada oferta, y es la misma implementación que sirve de línea base en la evaluación.
    \item \textbf{Fusión.} Las dos listas se combinan mediante fusión por rangos recíprocos (RRF) \citep{Cormack2009}: cada oferta suma $1/(k + r)$ por cada lista en la que aparece en la posición $r$, con $k = 60$. Como la fusión usa solo las posiciones, no es necesario calibrar entre sí la similitud coseno y el puntaje de BM25.
    \item \textbf{Reordenamiento.} Las 20 primeras ofertas de la lista fusionada pasan por un modelo de reordenamiento de tipo \textit{cross-encoder}, Qwen3-Reranker-0.6B \citep{Zhang2025Qwen3}, que lee juntas la consulta y la oferta \citep{Nogueira2019} y estima la probabilidad de que la oferta sea relevante. El modelo recibe el encabezado de la oferta (título, empresa, ubicación, tipo de contrato y salario) y sus fragmentos, hasta 2.000 caracteres. Se conservan las cinco ofertas con mayor probabilidad.
\end{enumerate}

El filtro se aplica dentro de cada motor, durante la búsqueda y no sobre sus resultados: Qdrant lo evalúa sobre la carga útil de los puntos, pgvector lo incorpora a la consulta SQL, con barridos iterativos cuando usa el índice, y BM25 descarta las ofertas que no lo cumplen antes de cortar en 50. Una oferta sin salario informado nunca cumple una condición de salario. Si ninguna oferta cumple el filtro, la búsqueda se repite sin él y el prompt indica al modelo que debe advertir al usuario que ninguna oferta satisface todas sus condiciones.

Dos mecanismos completan la etapa. El primero es una compuerta de \emph{sin resultados}: si incluso la mejor oferta reordenada tiene una probabilidad menor que 0,134, el sistema responde que no encontró ofertas relevantes en lugar de presentar resultados que no lo son. El umbral se calibró sobre un conjunto de consultas de desarrollo (Sección~\ref{sec:art-iter-recuperacion}). El segundo es la degradación ante fallos: si el servidor de reordenamiento no responde, se conserva el orden de la fusión.

\subsection{Construcción del contexto}

Para cada una de las cinco ofertas se arma un bloque de contexto con un encabezado (título, empresa, ubicación, tipo de contrato, salario y similitud) y el texto de sus fragmentos, limitado a 2.400 caracteres.

Antes de cerrar el contexto, una etapa de \emph{detección de campos} identifica si la consulta pregunta por salario, beneficios, requisitos, habilidades o responsabilidades. Si el campo solicitado no aparece en los fragmentos recuperados, se añade desde los datos estructurados de la oferta. Esta etapa reduce el riesgo de que el modelo complete con información plausible pero inexistente un dato que sí está en la base de datos y simplemente no fue recuperado.

\subsection{Generación}

El prompt de sistema instruye al modelo a responder \emph{únicamente} a partir de las ofertas proporcionadas, a indicar explícitamente cuando estas no contienen la información solicitada, a no inventar detalles y a responder en el idioma de la pregunta aunque las ofertas estén en otro. El historial reciente de la conversación, hasta 12 mensajes, se incluye para resolver preguntas de seguimiento; la conversación solo afecta a la generación y no a la recuperación, que usa la consulta actual. La generación usa el modelo \texttt{qwen/qwen3.8-27b} servido por Groq, con temperatura 0,2 en la configuración de referencia y respuesta incremental. En la ruta híbrida, los resultados de la recuperación y de la agregación se ejecutan en paralelo y se combinan en un único prompt; si una de las dos rutas falla, se responde con la otra.

\section{Iteración 1: corrección}
\label{sec:art-vectorial}

La primera iteración partió de una revisión del artefacto frente a su documentación, hecha al preparar la evaluación. Su objetivo era que lo medido correspondiera a lo que el artefacto dice hacer. Corrigió las fallas que la Sección~\ref{sec:art-fallas} clasifica, la mayoría en la ingesta, la agregación y la búsqueda vectorial, y rediseñó el almacenamiento vectorial para que pgvector y Qdrant pudieran compararse en igualdad de condiciones. Esta sección describe ese rediseño.

\subsection{Diagnóstico del almacenamiento vectorial}

En su versión inicial, JobAI almacenaba los vectores en una columna \texttt{vector(768)} de PostgreSQL, mediante la extensión pgvector \citep{Pgvector}. El esquema incluía un índice IVFFlat creado en la migración inicial, es decir, sobre una tabla todavía vacía y con los parámetros por defecto. La consulta de similitud, sin embargo, ordenaba los resultados por la expresión derivada $1 - d(\mathbf{q}, \mathbf{v})$ y filtraba por umbral sobre esa misma expresión. El planificador de PostgreSQL solo puede servir una búsqueda de vecinos cercanos desde el índice cuando la consulta ordena directamente por el operador de distancia (\texttt{ORDER BY embedding <=> q LIMIT k}).

La Tabla~\ref{tab:art-explain} resume la inspección con \texttt{EXPLAIN ANALYZE} sobre el corpus de evaluación (4.013 fragmentos), antes y después de la migración del índice.

\begin{table}[htbp!]
    \centering
    \caption{Plan de ejecución de la búsqueda de 15 vecinos en pgvector sobre el corpus de evaluación}
    \label{tab:art-explain}
    \begin{tabular}{|p{5.2cm}|p{2.2cm}|p{4.6cm}|r|}
        \hline
        \textbf{Consulta} & \textbf{Índice} & \textbf{Plan elegido} & \textbf{Tiempo} \\
        \hline
        Original (orden por $1-d$) & IVFFlat & Barrido secuencial + ordenamiento & 25,7 ms \\
        \hline
        Reescrita (orden por $d$) & IVFFlat & Recorrido del índice & 2,8 ms \\
        \hline
        Reescrita (orden por $d$) & HNSW & Barrido secuencial + ordenamiento & 24,8 ms \\
        \hline
        Reescrita, índice forzado & HNSW & Recorrido del índice & 2,0 ms \\
        \hline
    \end{tabular}
\end{table}

La inspección confirma tres hechos.
\begin{itemize}
    \item La consulta original nunca usaba el índice: recorría y ordenaba la tabla completa.
    \item El índice IVFFlat creado sobre la tabla vacía no perdía \textit{recall} (1,000 frente a la búsqueda exacta en 100 consultas), probablemente porque todos los vectores quedaron en una sola lista; en la práctica se comportaba como un recorrido completo.
    \item Con el índice HNSW y la consulta ya corregida, el planificador de PostgreSQL sigue prefiriendo el barrido secuencial a esta escala, aunque el recorrido del índice es 12 veces más rápido. El modelo de costos de pgvector sobrestima el costo de arranque del índice en tablas de pocos miles de filas. Por eso la búsqueda admite un modo que fuerza el índice, y la evaluación mide ambos comportamientos.
\end{itemize}

La misma inspección reveló que el umbral de similitud de 0,5 no filtra nada con el modelo de entonces, e5-base-v2. En 20 consultas de prueba, la similitud coseno mínima entre la consulta y \emph{cualquier} fragmento del corpus fue 0,66, y la mediana 0,75. Las similitudes de e5-base-v2 se concentran en valores altos, así que todos los fragmentos superan el umbral. La segunda iteración reemplazó este umbral por la compuerta de sin resultados descrita en la Sección~\ref{sec:art-recuperacion}.

Este diagnóstico cambia la motivación de la migración. La lentitud observada en pgvector no se debía al motor, sino a una consulta que no usaba el índice. Por eso la migración no se planteó como una mejora supuesta de antemano, sino como una alternativa que se pone a prueba: la hipótesis H3 (Sección~\ref{sec:hipotesis}) sostiene que, a la escala del corpus real, el motor no limita la calidad ni la latencia, y que las diferencias aparecen solo a mayor escala y con filtros selectivos.

\subsection{Diseño de la migración}

El diseño debía cumplir dos requisitos: que la plataforma pudiera servirse desde Qdrant \citep{QdrantDocs} o desde pgvector, y que ambos motores pudieran compararse en igualdad de condiciones. Las decisiones principales fueron cinco.

\begin{description}
    \item[Abstracción del almacén vectorial.] Se definió una interfaz única con estas operaciones:
    \begin{itemize}
        \item escritura de los fragmentos de una oferta;
        \item búsqueda con número de resultados, umbral, filtros estructurados, modo y parámetro $ef_{search}$;
        \item búsqueda restringida a un conjunto de ofertas;
        \item verificación de existencia y borrado.
    \end{itemize}
    Los módulos de ingesta, RAG y gestión de ofertas dependen de la interfaz y no de una implementación concreta. Una variable de configuración selecciona el motor que sirve a la aplicación.
    \item[Corrección de la consulta en pgvector.] La búsqueda se reescribió para ordenar por el operador de distancia, y el umbral se aplica después de obtener los $k$ vecinos, lo cual es equivalente porque la similitud decrece monótonamente con la distancia. El índice IVFFlat se reemplazó por un índice HNSW \citep{Malkov2020} con $m = 16$ y $ef_{construction} = 128$. La búsqueda admite tres modos:
    \begin{itemize}
        \item \emph{automático}: decide el planificador, que en tablas pequeñas puede preferir el barrido secuencial;
        \item \emph{aproximado}: fuerza el uso del índice;
        \item \emph{exacto}: fuerza la fuerza bruta, que es la referencia para medir el \textit{recall}.
    \end{itemize}
    Para búsquedas con filtros se activan los barridos iterativos del índice, mediante el parámetro \texttt{hnsw.\allowbreak iterative\_scan}. Sin ellos, un filtro selectivo puede devolver menos de $k$ resultados.
    \item[Colección en Qdrant.] La colección usa distancia coseno y los mismos parámetros HNSW que pgvector. Cada punto almacena como carga útil (\textit{payload}) el texto y el tipo del fragmento y los atributos de la oferta usados en los filtros: salario, tipo de contrato, modalidad y fuente, con índices de \textit{payload} sobre los campos filtrables. Los identificadores de fragmento de PostgreSQL no son UUID, por lo que el identificador de cada punto se deriva de forma determinista mediante un UUID versión 5. Se redujo además el umbral a partir del cual Qdrant construye el índice HNSW de un segmento: con el valor por defecto, un corpus de unos pocos miles de vectores se resolvería por fuerza bruta y la comparación no mediría el índice.
    \item[PostgreSQL como fuente de verdad.] El texto de los fragmentos y su relación con las ofertas permanecen en PostgreSQL, con borrado en cascada e integridad referencial. Al escribir, el almacén de Qdrant primero confirma la transacción en PostgreSQL y luego reemplaza los puntos de la oferta. La verificación de si una oferta ya tiene \textit{embeddings} se hace contra Qdrant, de modo que una escritura fallida en Qdrant se repara en el siguiente ciclo de ingesta.
    \item[Carga y verificación.] Un comando de línea copia en lotes todos los vectores existentes de PostgreSQL a Qdrant. La operación es idempotente, porque los identificadores son deterministas. Al terminar, el comando espera a que Qdrant indexe la colección y verifica que coincidan tres conteos: filas en PostgreSQL, puntos en Qdrant y vectores indexados.
\end{description}

Sobre el corpus de evaluación, con los vectores de e5-base-v2, la búsqueda exacta en ambos motores coincidió en el 99,85\,\% de los 10 primeros resultados de 200 consultas (99,75\,\% tras recodificar el corpus con el modelo multilingüe). Las diferencias corresponden a empates entre fragmentos de texto idéntico, que cada motor ordena de forma distinta. Esto confirma que los dos almacenes contienen los mismos datos y usan la misma métrica. La comparación de desempeño entre ambos motores se reporta en el Capítulo~\ref{cap:resultados}.

\section{Iteración 2: mejora}
\label{sec:art-iter-recuperacion}

\subsection{Diagnóstico}

Una primera corrida completa de la evaluación, con la configuración que resultó de la Iteración~1, y una revisión del pipeline frente a la literatura mostraron cuatro brechas entre lo que el trabajo mide y lo que el artefacto hacía. Las dos primeras motivan los componentes cuyo aporte se mide en el Capítulo~\ref{cap:resultados}; la tercera es una falla que la Sección~\ref{sec:art-fallas} clasifica (A6), y la cuarta motiva el cambio de modelo de \textit{embeddings}.
\begin{itemize}
    \item \textbf{Las restricciones de la consulta no llegaban a la búsqueda vectorial.} Una consulta con señales de recuperación, como \textit{\enquote{find remote Python jobs paying over 150k}}, iba directamente a la búsqueda semántica y sus condiciones se ignoraban. Cuando el clasificador sí extraía las condiciones, la plataforma las resolvía con un filtro SQL que exigía la coincidencia de subcadena del título (\texttt{ILIKE}), es decir, la misma búsqueda léxica que el trabajo pretende superar, y pasaba al modelo hasta 30 ofertas ordenadas por fecha y no por relevancia. Además, el filtro no contemplaba la modalidad remota. En consecuencia, el benchmark de búsqueda filtrada (PE4) medía una capacidad que la aplicación no usaba.
    \item \textbf{El pipeline correspondía al esquema más básico de RAG.} En la clasificación de \citet{Gao2023}, el pipeline era un RAG \emph{ingenuo}: indexación, una única búsqueda densa y generación. Carecía de las optimizaciones posteriores a la recuperación que caracterizan al RAG \emph{avanzado}, como el reordenamiento, y de la recuperación híbrida que combina representaciones dispersas y densas.
    \item \textbf{Las consultas en español recibían respuestas en inglés.} En la primera corrida, 21 de las 25 consultas en español se respondieron en inglés: sin una instrucción explícita, el modelo respondía en el idioma del contexto, que es el de las ofertas.
    \item \textbf{El modelo de \textit{embeddings} era solo para inglés.} e5-base-v2 \citep{Wang2022E5} se entrenó con texto en inglés, de modo que las consultas en español dependían de la coincidencia accidental de términos.
\end{itemize}

\subsection{Cambios}

La Tabla~\ref{tab:art-cambios} resume los cambios de la iteración. El diseño de cada uno se describió en la Sección~\ref{sec:art-consulta}; aquí se resume su motivación.

\begin{table}[htbp!]
    \centering
    \caption{Cambios de la iteración de recuperación}
    \label{tab:art-cambios}
    \small
    \begin{tabular}{|p{3cm}|p{5.4cm}|p{5.4cm}|}
        \hline
        \textbf{Aspecto} & \textbf{Antes} & \textbf{Después} \\
        \hline
        Restricciones de la consulta & Ignoradas si la consulta tenía señales de recuperación; si no, filtro SQL por subcadena del título, con 30 ofertas ordenadas por fecha & Extraídas por el clasificador (inglés y español) y aplicadas como filtro dentro de la búsqueda densa y de BM25 \\
        \hline
        Primera etapa & Búsqueda densa de 15 fragmentos con umbral de similitud 0,5 & Búsqueda densa de 50 fragmentos y BM25 de 50 ofertas, fusionadas con RRF \\
        \hline
        Orden final & Similitud del mejor fragmento & Reordenamiento con Qwen3-Reranker-0.6B sobre 20 candidatas \\
        \hline
        Consultas sin ofertas relevantes & Siempre se devolvían cinco ofertas & Compuerta sobre el mejor puntaje del reordenamiento (0,134) \\
        \hline
        Modelo de \textit{embeddings} & e5-base-v2 (inglés) & multilingual-e5-base \\
        \hline
        Idioma de la respuesta & El del contexto & El de la pregunta \\
        \hline
        Salarios fuera de rango & Se almacenaban tal como llegaban & Se registran como no informados \\
        \hline
    \end{tabular}
\end{table}

\begin{description}
    \item[Filtros en la recuperación.] Aplicar el filtro durante la búsqueda, y no sobre los $k$ vecinos ya recuperados, evita que un filtro selectivo deje la lista incompleta, el problema que motiva los métodos de búsqueda aproximada con filtros \citep{Gollapudi2023, Patel2024}. Que la búsqueda densa y la léxica usen el mismo filtro garantiza que la fusión no reintroduzca ofertas que no cumplen las condiciones.
    \item[Búsqueda híbrida.] BM25 aporta la coincidencia exacta de términos poco frecuentes, como nombres de tecnologías, que un modelo denso puede diluir. Se eligió RRF \citep{Cormack2009} porque no requiere calibrar puntajes de naturaleza distinta ni datos de entrenamiento, con el valor habitual $k = 60$ y sin ajustarlo. \citet{Bruch2023} encontraron, sin embargo, que RRF es sensible a ese parámetro y que una combinación convexa de los puntajes la supera, con pocos ejemplos para ajustar su único peso.
    \item[Reordenamiento.] Un \textit{cross-encoder} lee la consulta y la oferta en una misma pasada, por lo que puede juzgar la relación entre ambas con más detalle que la comparación de dos vectores calculados por separado \citep{Nogueira2019}. Su costo crece con el número de candidatas, por eso solo se aplica a las 20 primeras de la lista fusionada.
    \item[Modelo de \textit{embeddings} multilingüe.] multilingual-e5-base \citep{Wang2024mE5} conserva la dimensión de 768 de e5-base-v2, de modo que el cambio no alteró el esquema. Los 4.013 fragmentos del corpus de evaluación se volvieron a codificar en el lugar, con los mismos identificadores y textos. Los vectores de e5-base-v2 se conservaron en una colección aparte para medir el efecto del cambio de modelo (Capítulo~\ref{cap:resultados}).
    \item[Idioma de la respuesta y salarios.] Los tres prompts de generación incluyen ahora la instrucción de responder en el idioma de la pregunta. La regla de salarios plausibles se aplicó también al corpus congelado: afectó a 12 ofertas, entre ellas un salario de 86.528.000 dólares para un puesto de despacho y montos en coronas suecas.
\end{description}

\subsection{Selección del modelo de reordenamiento}

La elección del modelo de reordenamiento, del número de candidatas y del umbral de la compuerta son decisiones de configuración. Tomarlas mirando las consultas de prueba sesgaría la evaluación a favor del artefacto, así que se tomaron con un conjunto de desarrollo separado y con una regla de decisión fijada antes de ver los resultados. En términos de DSR, es una evaluación formativa: sus resultados orientan el diseño y no se reportan como resultados del artefacto.

El conjunto de desarrollo contiene 20 necesidades de información, cinco por categoría, formuladas en 28 consultas en inglés y en español, sin solapamiento con las de prueba, más 10 consultas fuera de dominio que solo se usan para el umbral. Los juicios de relevancia siguieron el mismo procedimiento que la evaluación (Sección~\ref{sec:met-juicios}). Se compararon dos modelos de reordenamiento, bge-reranker-v2-m3, basado en BGE-M3 \citep{Chen2024M3}, y Qwen3-Reranker-0.6B \citep{Zhang2025Qwen3}, en cinco variantes: con y sin instrucción de tarea, y con el orden final dado solo por el reordenamiento o fusionado por RRF con el orden de la primera etapa. La regla tenía tres etapas:
\begin{enumerate}
    \item Con 30 candidatas, se elige la variante de mayor nDCG@10. Si otra queda a 0,02 o menos, se prefiere la más simple.
    \item Con la variante elegida, se usan 20 candidatas en lugar de 30 si el nDCG@10 baja 0,01 o menos.
    \item El umbral es el mayor valor que deja por debajo, como máximo, a una de las 28 consultas de desarrollo, evaluado sobre el mejor puntaje de cada consulta. Se adopta solo si deja por debajo al menos a 8 de las 10 consultas fuera de dominio.
\end{enumerate}

\begin{table}[htbp!]
    \centering
    \caption{Selección del modelo de reordenamiento en el conjunto de desarrollo (28 consultas, 30 candidatas)}
    \label{tab:art-rerank-dev}
    \begin{tabular}{|l|c|c|}
        \hline
        \textbf{Variante} & \textbf{nDCG@10} & \textbf{P@5} \\
        \hline
        Sin reordenamiento (densa + BM25 + RRF) & 0{,}742 & 0{,}757 \\
        \hline
        bge-reranker-v2-m3 & 0{,}804 & 0{,}814 \\
        \hline
        bge-reranker-v2-m3, fusionado con RRF & 0{,}798 & 0{,}800 \\
        \hline
        Qwen3-Reranker-0.6B & 0{,}820 & 0{,}864 \\
        \hline
        Qwen3-Reranker-0.6B con instrucción & 0{,}839 & 0{,}879 \\
        \hline
        Qwen3-Reranker-0.6B con instrucción, fusionado con RRF & 0{,}812 & 0{,}879 \\
        \hline
    \end{tabular}
\end{table}

La Tabla~\ref{tab:art-rerank-dev} muestra la primera etapa. Todas las variantes mejoraron entre 0,06 y 0,10 el nDCG@10 de la recuperación sin reordenamiento. Qwen3-Reranker-0.6B sin instrucción quedó a 0,019 de la mejor variante, por lo que la regla lo eligió por ser más simple. En la segunda etapa, pasar de 30 a 20 candidatas cambió el nDCG@10 de 0,832 a 0,830 y redujo la mediana de la latencia de recuperación de 1.184 a 788~ms, así que se adoptaron 20. En la tercera, el mejor puntaje más bajo entre las consultas de desarrollo fue 0,132 y el siguiente 0,134; las diez consultas fuera de dominio quedaron por debajo de 0,09. El umbral de 0,134 cumplió la condición y se adoptó.

Una observación cualitativa motivó incluir la fusión con RRF entre las variantes: bge-reranker-v2-m3 es muy literal. Asignó una probabilidad de 0,97 al par \textit{\enquote{security engineer}} y \textit{\enquote{Security Engineer, Cloud}}, pero de 0,002 al par formado por la descripción funcional \textit{\enquote{Cybersecurity roles defending company systems against attacks}} y una oferta de \textit{\enquote{Security Engineer}}. Se temía que perjudicara las consultas de desajuste de vocabulario; en el conjunto de desarrollo no ocurrió, y la variante elegida no usa la fusión.

\section{Fallas detectadas por la evaluación formativa}
\label{sec:art-fallas}

En un sistema RAG muchos defectos son \emph{silenciosos}: el sistema sigue respondiendo con fluidez y de forma plausible, porque el modelo de lenguaje redacta una respuesta con lo que recibe, aunque lo recibido esté incompleto, mal filtrado o mal enrutado. Usando la aplicación, esos defectos no se notan. \citet{Barnett2024} concluyen que la validación de un sistema RAG solo es posible durante su operación. En JobAI, la preparación y la ejecución del protocolo de evaluación revelaron 17 fallas antes de esa etapa, 12 del artefacto y 5 del propio arnés de evaluación, además de dos comportamientos silenciosos del motor de almacenamiento que el diseño debe tener en cuenta. En términos de DSR son resultados de una evaluación \emph{formativa} \citep{Venable2016}: se obtuvieron durante el ciclo de diseño y orientaron sus correcciones.

Las fallas se clasifican en cuatro tipos según su efecto:
\begin{description}
    \item[A. Silenciosas que cambian los resultados.] El sistema devuelve ofertas, cifras o respuestas incorrectas sin ningún error visible.
    \item[B. Silenciosas que cambian el costo.] Los resultados son correctos, pero se obtienen con un costo mayor del necesario.
    \item[C. Visibles.] Rompen una funcionalidad o producen una respuesta evidentemente incorrecta.
    \item[D. Del arnés de evaluación.] No afectan al artefacto, sino a la medición, y habrían invalidado parte de la evaluación.
\end{description}
Los dos comportamientos del motor (M1 y M2) no se cuentan como fallas, porque no son defectos del artefacto; se describen aparte en la Sección~\ref{sec:art-motor}.
La Tabla~\ref{tab:art-fallas} las resume. Después se describe cada una con la misma secuencia: observación, causa, efecto, corrección y lección para otros sistemas RAG.

{\small
\begin{longtable}{|p{0.6cm}|p{4.4cm}|p{2.8cm}|p{3.6cm}|p{2.6cm}|}
    \caption{Fallas detectadas por la evaluación formativa. Las filas M son comportamientos del motor, no fallas del artefacto}
    \label{tab:art-fallas} \\
    \hline
    \textbf{Id} & \textbf{Falla} & \textbf{Detección} & \textbf{Efecto medido} & \textbf{Estado} \\
    \hline
    \endfirsthead
    \hline
    \textbf{Id} & \textbf{Falla} & \textbf{Detección} & \textbf{Efecto medido} & \textbf{Estado} \\
    \hline
    \endhead
    A1 & La personalización por currículo reemplazaba la consulta & Revisión frente a la documentación & Mismas ofertas para cualquier pregunta & Corregida (It.~1) \\
    \hline
    A2 & Las plantillas de \enquote{remoto} buscaban la palabra en la ubicación & Revisión de la agregación & 44 ofertas reconocidas frente a 682 & Corregida (It.~1) \\
    \hline
    A3 & El umbral de similitud de 0,5 no filtraba & Inspección de similitudes & Mínima 0,66 y mediana 0,75 en 20 consultas & Reemplazado por compuerta (It.~2) \\
    \hline
    A4 & Las restricciones de la consulta no llegaban a la búsqueda & Primera corrida & C@5 de 0,41 sin filtros frente a 0,95 con filtros & Corregida (It.~2) \\
    \hline
    A5 & El normalizador no convierte monedas ni montos por hora & Respuestas de agregación & Salario promedio del corpus de 544.127 USD & Mitigada (It.~2) \\
    \hline
    A6 & Respuestas en inglés a consultas en español & Primera corrida (PE2) & 21 de 25 respuestas & Corregida (It.~2): 4 de 24 \\
    \hline
    A7 & La plantilla de ofertas recientes se atendía por la ruta semántica & Corrida de prueba & Exactitud de la agregación 0,667 & Corregida: 0,733 \\
    \hline
    A8 & La palabra \textit{hiring} activa el patrón de recuperación & Corrida de prueba & 2 de 15 consultas de agregación mal enrutadas & Abierta \\
    \hline
    B1 & La consulta ordenaba por $1-d$ y no usaba el índice & \texttt{EXPLAIN ANALYZE} & 25,7~ms frente a 2,8~ms & Corregida (It.~1) \\
    \hline
    B2 & Índice IVFFlat creado sobre la tabla vacía & \texttt{EXPLAIN ANALYZE} & Equivalía a un recorrido completo & Reemplazado por HNSW (It.~1) \\
    \hline
    C1 & La ruta de agregación fallaba sin plantilla & Revisión del código & Error de ejecución & Corregida (It.~1) \\
    \hline
    C2 & Faltan plantillas: contar con una condición y contar empresas distintas & Corrida de prueba & Respuesta incorrecta con enrutamiento correcto & Abierta \\
    \hline
    D1 & El juez sin condiciones duras aceptaba ofertas que no cumplían la restricción & Revisión de juicios & La métrica no distinguía si un sistema respeta las restricciones & Corregida: juez v2 \\
    \hline
    D2 & Ventana de 4.096 \textit{tokens} del juez de RAGAS & Revisión de salidas del juez & Fidelidad truncada en 21 de 85 respuestas & Corregida: 16.384 \\
    \hline
    D3 & El juez local asigna fidelidad 0 a cifras presentes en el contexto & Revisión de casos extremos & PE2 no confiable con el juez local & Pendiente: juez final \\
    \hline
    D4 & Exportación limitada a 20 registros e ingesta no desactivable & Preparación del corpus & Habrían alterado el corpus congelado & Corregidas (It.~1) \\
    \hline
    D5 & El cálculo del kappa usaba la versión anterior del prompt del juez & Revisión del arnés & Habría comparado 269 de los 401 pares con veredictos del prompt anterior & Corregida \\
    \hline
    M1 & Un filtro selectivo en pgvector sin barrido iterativo devuelve listas incompletas & Benchmark (PE4) & Corpus real, filtro del 50\,\%: listas completas en el 22\,\% de las consultas & Evitado en el diseño (It.~1) \\
    \hline
    M2 & Con HNSW y pocas filas, el planificador prefiere el barrido secuencial & \texttt{EXPLAIN ANALYZE} & 24,8~ms frente a 2,0~ms con el índice forzado & Documentada; modo que fuerza el índice \\
    \hline
\end{longtable}
}

\subsection{Fallas silenciosas que cambian los resultados}

\begin{description}
    \item[A1. Personalización por currículo.] La documentación indicaba que el perfil del usuario se antepone a su consulta; la implementación construía la búsqueda solo con el perfil. Un usuario con currículo recibía las mismas ofertas fuera cual fuera su pregunta, y el modelo redactaba para cada pregunta una respuesta coherente con esas ofertas. Se corrigió para anteponer el perfil. \emph{Lección:} una ruta de personalización debe probarse variando la consulta y verificando que el resultado cambie.
    \item[A2. Ofertas remotas en la agregación.] Las plantillas que cuentan y listan ofertas remotas buscaban la palabra \textit{remote} en el texto de la ubicación y no usaban el atributo de modalidad que la ingesta obtiene de la fuente o del análisis estructurado. En el corpus de evaluación reconocían 44 ofertas remotas, frente a 682 marcadas con el atributo. El modelo informaba la cifra sin advertir nada. Se corrigió para usar el atributo, con el texto de la ubicación como respaldo. \emph{Lección:} cada plantilla de agregación debe contrastarse con una consulta de referencia escrita a mano, porque el modelo presenta cualquier cifra con la misma seguridad.
    \item[A3. Umbral de similitud.] El pipeline descartaba los fragmentos con similitud menor que 0,5. Con e5-base-v2, la similitud mínima entre la consulta y cualquier fragmento fue 0,66 en 20 consultas, y la mediana 0,75, de modo que el umbral nunca descartaba nada y toda consulta, incluso fuera de dominio, recibía cinco ofertas. La Iteración~2 lo reemplazó por una compuerta calibrada sobre el puntaje del reordenamiento (Sección~\ref{sec:art-iter-recuperacion}). \emph{Lección:} los umbrales de similitud dependen de la distribución de cada modelo y deben calibrarse con consultas dentro y fuera de dominio.
    \item[A4. Restricciones ignoradas.] Las consultas con señales de recuperación iban directamente a la búsqueda semántica y sus condiciones de salario, modalidad o tipo de contrato se ignoraban (Sección~\ref{sec:art-iter-recuperacion}). Las respuestas mencionaban ofertas que no cumplían lo pedido, sin advertirlo. Sin filtros, solo el 41\,\% de las cinco primeras ofertas de las consultas con restricciones las cumplía; con los filtros de la Iteración~2, el 95\,\%. \emph{Lección:} la calidad de la recuperación debe medirse también con criterios deterministas, como el cumplimiento de las restricciones, que no dependen del juez.
    \item[A5. Normalización de salarios.] El normalizador no convierte monedas ni algunos montos por hora o mensuales. Entre otros casos, una oferta de despacho registraba 86.528.000 dólares anuales, y el salario promedio del corpus que respondía la ruta de agregación era de 544.127 dólares. La Iteración~2 registró como no informados los salarios fuera del intervalo de 10.000 a 1.000.000 dólares anuales o con un mínimo mayor que el máximo; la regla afectó a 12 ofertas y el promedio pasó a 158.786 dólares. Es una mitigación: el salario de esas ofertas se pierde en lugar de convertirse. \emph{Lección:} las respuestas de agregación heredan los errores de normalización de los datos, y deben revisarse los valores extremos antes de evaluar.
    \item[A6. Idioma de la respuesta.] En la primera corrida, 21 de las 25 consultas en español se respondieron en inglés: sin una instrucción explícita, el modelo respondía en el idioma del contexto. La Iteración~2 agregó a los tres prompts la instrucción de responder en el idioma de la pregunta; 4 de 24 respuestas siguieron en inglés. \emph{Lección:} el idioma de la respuesta es una propiedad que debe especificarse y medirse, no suponerse.
    \item[A7. Ofertas recientes.] Ante \textit{\enquote{what are the most recent job postings}}, el clasificador devolvía la etiqueta de recuperación junto con la plantilla de ofertas recientes, y la consulta se atendía por la ruta semántica con una respuesta plausible pero no ordenada por fecha. Se corrigió con una regla de enrutamiento: las plantillas de estadísticas, valores distintos y ofertas recientes se atienden siempre por la agregación (Sección~\ref{sec:art-clasificacion}). El Capítulo~\ref{cap:resultados} reporta la exactitud antes y después. \emph{Lección:} la salida del clasificador debe validarse con reglas deterministas cuando la etiqueta y la plantilla se contradicen.
    \item[A8. La palabra \textit{hiring}.] El patrón de recuperación incluye la palabra \textit{hiring}, de modo que \textit{\enquote{list all companies hiring engineers}} y \textit{\enquote{how many different companies are hiring}} se enrutan como consultas mixtas. La falla se detectó en el conjunto de prueba y no se corrigió: corregirla habría ajustado el artefacto a esas consultas. Queda como trabajo futuro. \emph{Lección:} los patrones léxicos de enrutamiento deben probarse con consultas que usan las mismas palabras con otra intención.
\end{description}

\subsection{Fallas silenciosas que cambian el costo}

\begin{description}
    \item[B1. Consulta que no usaba el índice.] La búsqueda ordenaba por la expresión $1 - d(\mathbf{q}, \mathbf{v})$ en lugar de por la distancia, y el planificador de PostgreSQL no puede servir ese orden desde el índice. Los resultados eran correctos, porque el barrido completo es exacto, pero la búsqueda tardaba 25,7~ms en lugar de 2,8~ms (Tabla~\ref{tab:art-explain}). Se corrigió en la Iteración~1. \emph{Lección:} el plan de ejecución de cada consulta vectorial debe inspeccionarse; un índice que existe no garantiza que se use.
    \item[B2. Índice creado sobre la tabla vacía.] El índice IVFFlat se creó en la migración inicial, antes de cargar datos, de modo que todos los vectores quedaron en una sola lista y el índice equivalía a un recorrido completo. Se reemplazó por HNSW, que no requiere entrenamiento. \emph{Lección:} los índices que se entrenan sobre los datos deben construirse después de cargarlos.
\end{description}

\subsection{Fallas visibles}

\begin{description}
    \item[C1. Agregación sin plantilla.] Cuando los patrones clasificaban una consulta como agregación sin seleccionar ninguna plantilla, la ejecución fallaba con un error. Las pruebas de integración no lo detectaban porque simulaban el repositorio SQL. Se corrigió delegando la selección de la plantilla al clasificador basado en LLM, con la recuperación como alternativa, y las pruebas verifican ahora la plantilla seleccionada. \emph{Lección:} las pruebas que simulan el acceso a datos pueden ocultar fallas de integración.
    \item[C2. Plantillas faltantes.] El catálogo no tiene una plantilla que cuente las ofertas que cumplen una condición de salario, ni una que cuente empresas distintas. Ante \textit{\enquote{how many jobs pay more than 200k}}, el sistema elige correctamente la plantilla más cercana, que lista ofertas, y la respuesta no contiene el conteo pedido. Se detectó en el conjunto de prueba y queda abierta, por la misma razón que A8. \emph{Lección:} un catálogo cerrado de plantillas debe evaluarse por cobertura de las preguntas, no solo por exactitud del enrutamiento.
\end{description}

\subsection{Fallas del arnés de evaluación}

\begin{description}
    \item[D1. Juez sin condiciones duras.] La primera versión del prompt del juez calificaba como relevante una oferta del cargo correcto aunque no cumpliera el salario, la modalidad o el tipo de contrato pedidos. Con ese juez, la métrica no distinguía a un sistema que respeta las restricciones de uno que las ignora. La segunda versión trata esas condiciones como duras (Sección~\ref{sec:met-juicios}). Este cambio, junto con el del agrupamiento, explica que el nDCG@10 de la recuperación densa con e5-base-v2 pasara de 0,73 en la primera corrida a 0,68 en la definitiva sin que cambiara el sistema (Capítulo~\ref{cap:resultados}). \emph{Lección:} el prompt del juez debe expresar las condiciones que la métrica pretende medir.
    \item[D2. Ventana de contexto del juez de RAGAS.] Con la ventana por defecto de 4.096 \textit{tokens}, el contexto recuperado y las instrucciones no cabían, y la verificación de la fidelidad quedó truncada en 21 de las 85 respuestas. Todas se puntuaron de nuevo con una ventana de 16.384. \emph{Lección:} el arnés debe verificar que el juez recibe completo lo que juzga.
    \item[D3. Errores del juez local.] El juez local asignó fidelidad 0 a respuestas cuyas cifras y títulos aparecen literalmente en el contexto. Por eso la fidelidad medida con ese juez no se considera fiable, y se repetirá con el juez final y una verificación humana. \emph{Lección:} los puntajes extremos de un juez automático deben revisarse a mano antes de interpretarse.
    \item[D4. Exportación e ingesta.] El listado usado para exportar el catálogo estaba limitado a 20 registros, y la ingesta programada no podía desactivarse. Lo primero habría truncado el corpus exportado y lo segundo lo habría modificado durante la evaluación. Ambos se corrigieron antes de congelar el corpus. \emph{Lección:} la reproducibilidad exige poder congelar los datos y verificar su integridad.
    \item[D5. Versión del prompt en el cálculo del kappa.] El script que calcula el acuerdo entre el juez y los juicios humanos seleccionaba los veredictos de la primera versión del prompt, fijada en su código, mientras que la muestra se extrajo de los juicios de la segunda. Como la caché conserva los veredictos de ambas versiones, el cálculo no habría fallado: habría emparejado solo 269 de los 401 pares, con los veredictos del prompt anterior. Se detectó al preparar la elección del juez final, antes de calcular ningún kappa, y se corrigió haciendo explícita la versión del prompt. \emph{Lección:} cada artefacto del arnés que filtra resultados debe verificar que la cantidad de elementos emparejados es la esperada.
\end{description}

\subsection{Comportamientos del motor}
\label{sec:art-motor}

La evaluación reveló además dos comportamientos de pgvector que no son defectos del artefacto, pero que tienen la misma naturaleza silenciosa: el motor devuelve resultados sin ningún aviso, aunque estén incompletos o se obtengan por un camino más lento. El diseño debe tenerlos en cuenta.

\begin{description}
    \item[M1. Filtros sin barrido iterativo en pgvector.] Con el índice HNSW, pgvector aplica el filtro sobre los candidatos que recorre el índice; si el filtro es selectivo, puede devolver menos de $k$ resultados sin ningún aviso. En el corpus real, con un filtro del 50\,\% y $ef_{search}=16$, solo el 22\,\% de las consultas obtuvo la lista completa y el \textit{recall} fue de 0,74; con barridos iterativos, todas la obtuvieron y el \textit{recall} subió a 0,93. El artefacto nunca lo sufrió, porque antes de la Iteración~2 los filtros no llegaban a la búsqueda vectorial, y el diseño de la Iteración~1 ya activaba esos barridos en toda búsqueda filtrada. \emph{Lección:} una búsqueda filtrada debe verificar que devuelve el número de resultados pedido.
    \item[M2. Preferencia del planificador.] Aun con la consulta corregida y un índice HNSW, el planificador de PostgreSQL elige el barrido secuencial a la escala del corpus real: 24,8~ms frente a 2,0~ms con el índice forzado. No es un error del artefacto sino un comportamiento del motor; la búsqueda admite un modo que fuerza el índice y la evaluación mide ambos. \emph{Lección:} la latencia de un motor vectorial depende de las decisiones de su planificador, que deben medirse a la escala real de uso.
\end{description}

\subsection{Síntesis}

De las 12 fallas del artefacto, 10 son silenciosas: ocho cambian los resultados y dos el costo. Ninguna producía un error visible ni una respuesta evidentemente incorrecta, porque el modelo de lenguaje convertía en texto fluido lo que recibía. Se detectaron porque el protocolo mide salidas intermedias y no solo la respuesta final: el plan de ejecución de la consulta, el número de resultados de una búsqueda filtrada, el cumplimiento de las restricciones, la plantilla elegida, el idioma de la respuesta y los veredictos del juez. El Capítulo~\ref{cap:resultados} retoma esta observación en la discusión.

    
    % Capitulo 5
    \chapter{Resultados y evaluaci\'on}
    \label{cap:resultados}
    \pendiente{Este capítulo se redacta con los resultados de la evaluación (semanas del 19-oct al 1-nov). La estructura prevista es la siguiente.}

\section{Corpus y conjunto de consultas}
\pendiente{Estadísticas del corpus congelado: ofertas por fuente, fragmentos por tipo, tasa de análisis estructurado fallido. Descripción del conjunto de consultas por categoría e idioma.}

\section{Validación del juez de relevancia}
\pendiente{Tabla: kappa binario y ponderado (con IC 95\,\%) del juez local y del juez final frente a los juicios humanos; matriz de confusión.}

\section{Calidad de recuperación (PE1)}
\pendiente{Tabla: P@5, P@10, nDCG@10, MRR y Recall@10 por sistema (ILIKE, BM25, pgvector exacto, pgvector HNSW, Qdrant HNSW), con significancia frente a cada línea base. Figura: nDCG@10 por categoría de consulta. Tabla: resultados del subconjunto en español. Exactitud del enrutamiento en consultas de agregación.}

\section{Fidelidad de la generación (PE2)}
\pendiente{Tabla: fidelidad y relevancia de la respuesta por categoría; análisis de casos de baja fidelidad.}

\section{Latencia del pipeline (PE3)}
\pendiente{Tabla: p50 y p95 por etapa (\textit{embedding}, búsqueda, construcción del contexto, TTFT, total) con cada almacén vectorial.}

\section{pgvector frente a Qdrant (PE4)}

\subsection{Resultado preliminar sobre el corpus real}

La Tabla~\ref{tab:res-ann-real} resume una primera comparación sobre el corpus de evaluación (4.013 fragmentos). Se usaron 200 consultas formadas por títulos de ofertas del corpus y 3 repeticiones; la referencia es la búsqueda exacta en pgvector. La latencia se mide desde el cliente y corresponde a la búsqueda de 10 vecinos, sin el \textit{embedding} de la consulta.

\begin{table}[htbp!]
    \centering
    \caption{Recall@10 y latencia de búsqueda sobre el corpus real (resultado preliminar)}
    \label{tab:res-ann-real}
    \begin{tabular}{|l|c|c|c|}
        \hline
        \textbf{Configuración} & \textbf{Recall@10} & \textbf{p50 (ms)} & \textbf{p95 (ms)} \\
        \hline
        pgvector, exacta & 1{,}000 & 24{,}3 & 27{,}2 \\
        \hline
        Qdrant, exacta & 0{,}999 & 3{,}9 & 4{,}8 \\
        \hline
        pgvector, plan automático & 1{,}000 & 23{,}6 & 25{,}3 \\
        \hline
        pgvector, HNSW $ef=16$ & 0{,}976 & 2{,}0 & 2{,}7 \\
        \hline
        Qdrant, HNSW $ef=16$ & 0{,}925 & 3{,}0 & 3{,}7 \\
        \hline
        pgvector, HNSW $ef=64$ & 0{,}995 & 2{,}2 & 2{,}8 \\
        \hline
        Qdrant, HNSW $ef=64$ & 0{,}988 & 3{,}1 & 3{,}6 \\
        \hline
        pgvector, HNSW $ef=256$ & 0{,}999 & 3{,}2 & 3{,}8 \\
        \hline
        Qdrant, HNSW $ef=256$ & 0{,}998 & 3{,}5 & 4{,}1 \\
        \hline
    \end{tabular}
\end{table}

A esta escala se observan tres hechos que el benchmark completo deberá contrastar:
\begin{itemize}
    \item Con el plan automático, pgvector recorre la tabla completa (Sección~\ref{sec:art-vectorial}), por lo que la aplicación servida desde pgvector tarda unas ocho veces más que servida desde Qdrant.
    \item Forzando el índice, pgvector es ligeramente más rápido que Qdrant y alcanza un \textit{recall} mayor con el mismo $ef_{search}$. A este tamaño, la latencia de Qdrant está dominada por la sobrecarga del cliente y del protocolo.
    \item Las diferencias de \textit{recall} entre los dos motores en modo exacto se deben a empates entre fragmentos idénticos.
\end{itemize}

\pendiente{Figuras: curvas \textit{recall}@10 frente a latencia por escala (corpus real, 100k, 500k) y por selectividad del filtro (sin filtro, 50\,\%, 10\,\%, 1\,\%). Tabla: tiempo de construcción, tamaño del índice, memoria pico y QPS.}

\section{Discusión}
\pendiente{Respuesta a cada pregunta de evaluación, relación con la literatura revisada en el Capítulo~\ref{cap:estado-arte} y limitaciones observadas.}

    
    % Capitulo 10
    \chapter{Conclusiones y recomendaciones}
    \label{cap:conclusiones}
    \section{Conclusiones}

\pendiente{Una conclusión por objetivo específico, sustentada en los resultados del Capítulo~\ref{cap:resultados}:}
\begin{enumerate}
    \item \pendiente{Limitaciones de la búsqueda léxica y componentes de RAG que la literatura propone (objetivo 1).}
    \item \pendiente{Protocolo de evaluación reproducible y validez del juez (objetivo 2).}
    \item \pendiente{Las dos iteraciones y las fallas silenciosas que el protocolo reveló (objetivo 3).}
    \item \pendiente{Aporte y costo de cada componente frente a BM25, y efecto del motor vectorial (objetivo 4; hipótesis H1 a H3).}
\end{enumerate}

\section{Recomendaciones y trabajo futuro}

\pendiente{Líneas previstas:}
\begin{itemize}
    \item \pendiente{Evaluación con usuarios reales.}
    \item \pendiente{Diálogo de varios turnos, con el historial también en la recuperación.}
    \item \pendiente{Personalización por currículo, evaluada aparte.}
    \item \pendiente{Fusión híbrida ajustada (combinación convexa o RRF con $k$ calibrado en desarrollo) y BM25 con traducción automática de la consulta.}
    \item \pendiente{Reordenamiento más barato o condicional.}
    \item \pendiente{Fallas abiertas: patrón de \textit{hiring}, plantillas faltantes y extracción del tipo de contrato \enquote{prácticas}.}
    \item \pendiente{Normalización de monedas y periodos de salario.}
\end{itemize}

    
    
    % Anexos (Opcional)
    % \chapter*{Anexos}
    % % \section{Manual de instalación}

% \section{Manual de usuario}

% \section{Informe de ejecución de pruebas}


    
    % Indices (Opcional)
    % \chapter*{\'Indices}
    % \input{BackMatter/B10-Index}
    
	% Incluye las referencias
    \bibliographystyle{apalike}
 	\bibliography{Bibliography/referencias}

\end{document}
